\documentclass[12pt, oneside]{article}

\usepackage{amsmath, amsthm, amssymb, amsfonts}
\usepackage{mathtools}
\usepackage{geometry}
\usepackage{hyperref}
\usepackage{booktabs}
\usepackage{array}
\usepackage{xcolor}
\usepackage{titlesec}
\usepackage{abstract}
\usepackage{microtype}
\usepackage{setspace}
\usepackage[numbers, sort&compress]{natbib}

\geometry{
  paper=a4paper,
  top=2.5cm, bottom=2.5cm,
  left=3cm, right=3cm
}

\onehalfspacing

\hypersetup{
  colorlinks=true,
  linkcolor=black,
  citecolor=black,
  urlcolor=black
}

% Theorem environments
\newtheorem{theorem}{Theorem}[section]
\newtheorem{lemma}[theorem]{Lemma}
\newtheorem{proposition}[theorem]{Proposition}
\newtheorem{corollary}[theorem]{Corollary}
\newtheorem{definition}[theorem]{Definition}
\newtheorem{assumption}[theorem]{Assumption}
\newtheorem{remark}[theorem]{Remark}

% Preliminary section styling
\newcommand{\preliminary}[1]{%
  \colorbox{yellow!20}{\parbox{\dimexpr\linewidth-2\fboxsep}{%
    \textbf{\textcolor{red!70!black}{[PRELIMINARY OUTLINE --- NOT YET DRAFTED]}}\\[4pt]
    #1
  }}%
}

% Math operators
\DeclareMathOperator{\supp}{supp}
\DeclareMathOperator{\vol}{Vol}
\DeclareMathOperator{\interior}{int}
\DeclareMathOperator{\diam}{diam}
\newcommand{\norm}[1]{\left\|#1\right\|}
\newcommand{\abs}[1]{\left|#1\right|}
\newcommand{\inner}[2]{\langle #1, #2 \rangle}
\newcommand{\Xmod}{\mathcal{X}/{\sim}}
\newcommand{\Xquot}{\mathcal{X}/G}
\newcommand{\Ftil}{\widetilde{\Pi}}

\title{%
  \Large\textbf{Yarncrawler: World-State Reconstruction}\\[6pt]
  \large\textbf{as Constraint Closure over a Generative Field}
}

\author{Flyxion}
\date{}

\begin{document}

\maketitle

\begin{abstract}
This paper introduces Yarncrawler as a framework for reconstructing physical world-states
from multi-modal sensor observations under continuous dynamical constraint. The central
object is a consistency operator $\mathcal{C}$ acting on the quotient configuration space
$\mathcal{X}/{\sim}$ of coupled Relativistic Scalar-Vector Plenum field triples
$(\Phi, \mathbf{v}, S)$. Each sensor modality defines a projection operator
$\Pi_i : \mathcal{X} \to \mathcal{Y}_i$; each observation $y_i$ constrains the set of
field configurations consistent with what the sensor reports. The consistency operator
searches for the unique admissible configuration whose forward projections are compatible
with all observations simultaneously.

The main result---the Yarncrawler Identifiability Theorem---establishes that $\mathcal{C}$
has a unique fixed point if and only if the induced map
$\widetilde{\Pi} : \Xmod \to \prod_i \mathcal{Y}_i$ is injective on the feasible set
$\mathcal{F} = \Omega_{\mathrm{obs}} \cap \Gamma$. The Dynamical Disambiguation Theorem
extends this to realistic sensing regimes where static observations underdetermine the
world-state: trajectory fragments constrain the tangent bundle of the admissible manifold,
removing degrees of freedom that no static projection can touch.

The paper further establishes that the variational consistency framework and the
sheaf-theoretic repair framework---in which sensor modalities define local charts, observations
define local sections, and closure corresponds to the existence of a unique global
section---are formally equivalent. There exists a functor
$\Psi : \mathbf{Proj}(\mathcal{X}) \to \mathbf{Sh}(\Xmod)$ under which the consistency
operator is the variational realization of the sheaf repair operator, projection degeneracy
corresponds to non-trivial $H^1$, and entropy collapse corresponds to cohomology
trivialization. The RSVP entropy field $S$ is identified as the local density of
cohomological obstruction; its minimization under $\mathcal{C}$ is the physical expression
of sheaf repair.

Deployment at planetary, regional, and site scales is unified by treating each scale as a
different geometry of the feasible set under varying noise profiles. The framework is
implemented as a sensor fusion architecture combining learned trajectory operators with
physics-constrained forward models, applied to coupled ecosystem and civilization dynamics
calibrated from before-and-after observational pairs across modalities including infrared,
LiDAR, RGB, audio, and photonic strip sensors.
\end{abstract}

\tableofcontents
\newpage

%% ============================================================
\section{Introduction}
%% ============================================================

\subsection{The Problem}

Given a family of projection operators $\{\Pi_i\}_{i \in \mathcal{I}}$ mapping an unknown
world-state $[X] \in \Xmod$ to observable quantities $y_i \in \mathcal{Y}_i$, and given a
dynamical law governing how world-states evolve, the central question of this paper is:

\begin{quote}
\textit{Under what conditions does there exist a unique trajectory in the quotient
configuration space $\Xmod$ consistent with all projections simultaneously and admissible
under the dynamics?}
\end{quote}

This is not a modeling question. It is an identifiability question. The paper does not
propose a method for approximating world-states from sensor data---it establishes the
conditions under which a unique world-state is \emph{demanded} by the data, and proves
that when those conditions hold, a specific operator converges to it.

The setting is the following. The world-state is a coupled field triple
$X = (\Phi, \mathbf{v}, S)$---scalar density, vector flow, and entropy
production---evolving over a compact spatial domain under the Relativistic Scalar-Vector
Plenum (RSVP) field equations. The projections $\Pi_i$ correspond to sensor modalities:
infrared, LiDAR, RGB, audio, photonic strips. Each observation $y_i$ constrains which
field configurations are consistent with what the sensor reports. The consistency operator
$\mathcal{C}$ searches for the configuration that minimizes total projection residual
subject to physical admissibility. The dynamical admissibility set $\Gamma$ restricts
attention to configurations that lie on trajectories the world could have taken. The
feasible set $\mathcal{F} = \Omega_{\mathrm{obs}} \cap \Gamma$ is the intersection of all
constraints. Closure---the condition that $|\mathcal{F}| = 1$---is the condition that the
world-state is uniquely determined.

\subsection{The Central Result}

The paper proves the following.

The problem of multi-modal consistency---finding a world-state whose forward projections
are compatible with all available observations---is equivalent to the existence of a unique
global section of a sheaf $\mathcal{S}_\Pi$ induced by the projection family over a cover
of $\Xmod$.

More precisely: the consistency operator $\mathcal{C}$ is the variational realization of
the sheaf repair operator. Its fixed points are global sections. Its failure modes---
projection degeneracy and regularizer flatness---are cohomological obstructions $H^1 \ne 0$
and degenerate $H^0$. Adding sensors refines the cover and kills $H^1$ classes.
Strengthening the dynamical constraint extends the cover into the temporal direction and
kills residual ambiguity that spatial sensing cannot reach. Closure is the condition
$H^0 = 1$, $H^1 = 0$: the sheaf has exactly one global section.

This is the Yarncrawler Identifiability Theorem and its dynamical extension, proved in
Sections \ref{sec:identifiability} and \ref{sec:disambiguation} and their appendices.

\subsection{The Duality}

The framework has two faces, both present in this paper.

The \emph{forward} direction takes a world-state as given and asks how a system maintains
coherence under perturbation: how it detects tears in its semantic fabric, applies repair
morphisms that trivialize \v{C}ech cocycles, exports entropy through stigmergic
reinforcement, and sustains homeorhetic flows rather than static equilibria.

The \emph{inverse} direction takes sensor observations as given and asks what world-state
they structurally demand. Given what the sensors report, what field configuration is forced?

These are not complementary perspectives on the same phenomenon. They are dual realizations
of the same mathematical object. Both directions study the sheaf $\mathcal{S}_\Pi$ over the
quotient configuration space $\Xmod$. The forward direction studies the internal dynamics
of sections within $\mathcal{S}_\Pi$---how they evolve, heal, and accumulate into global
coherence. The inverse direction studies the conditions under which a unique global section
exists and the operator that converges to it. The RSVP field triple $(\Phi, \mathbf{v}, S)$
parametrizes the stalks in both directions. The Yarncrawler framework is the unified theory
of that sheaf and its dynamics.

The object common to both is neither the operator nor the sheaf taken separately. It is
the equivalence between them---the functor
$\Psi : \mathbf{Proj}(\mathcal{X}) \to \mathbf{Sh}(\Xmod)$ proved in
Section~\ref{sec:equivalence}, which shows that a projection family and a sheaf over a
configuration space are the same structure viewed from outside and inside respectively.

\subsection{The Consequence}

When closure is achieved---when $|\mathcal{F}| = 1$---three things happen simultaneously
that were previously three separate conditions.

The projection residuals collapse to within noise tolerance:
$\|\Pi_i(X^*) - y_i\| \le \epsilon_i$ for all $i$. The cohomology of the induced sheaf
trivializes: $H^1(\mathcal{U}, \mathcal{S}_\Pi) = 0$ and $\dim H^0 = 1$. The RSVP entropy
field $S$ is minimized subject to admissibility: no further constraint can be satisfied
that is not already satisfied.

These are not three criteria that must be checked independently. Under the
Sheaf-Variational Equivalence of Section~\ref{sec:equivalence}, they are one criterion
expressed in three languages---variational, cohomological, and field-theoretic. The entropy
of the RSVP field at the fixed point is the local density of cohomological obstruction,
and its minimization is the physical expression of sheaf repair.

The closed state $[X^*]$ is therefore not a best estimate. It is the unique trajectory
consistent with all projections and admissible dynamics---the world-state that no
alternative history can explain away.

\subsection{Structure of the Paper}

Section~\ref{sec:equivalence} establishes the Sheaf-Variational Equivalence. Sections
\ref{sec:worldstate} and \ref{sec:projections} develop the prerequisite machinery.
Section~\ref{sec:feasible} defines the feasible set geometrically.
Sections~\ref{sec:identifiability} and \ref{sec:disambiguation} contain the core theorems.
Section~\ref{sec:inversion} develops the inversion layer architecture.
Section~\ref{sec:closure} defines closure operationally and proves equivalence of the three
closure criteria. Section~\ref{sec:deployment} unifies planetary, regional, and site-scale
deployment. Section~\ref{sec:related} positions the framework against prior work.
Section~\ref{sec:epistemology} draws the epistemological consequence. Appendices
\ref{app:frechet}--\ref{app:operators} provide the functional-analytic foundations.

%% ============================================================
\section{The Sheaf-Variational Equivalence}
\label{sec:equivalence}
%% ============================================================

\subsection{Two Descriptions of the Same Object}

The Yarncrawler framework has two faces. The first describes a self-refactoring system
that maintains coherence by repairing tears in its semantic fabric: local sections are
parsed, inconsistencies are detected as \v{C}ech cocycles, and repair morphisms trivialize
them until a global section exists. The second describes a consistency operator that
searches for a field configuration whose forward projections are compatible with all
available sensor observations simultaneously. These descriptions appear to operate at
different levels of abstraction---one categorical and topological, the other variational
and analytic. This section proves they are the same description.

The identification is not a metaphor. There is a precise functor between the category of
sheaves of semantic sections over a cover and the category of projection families on the
quotient configuration space $\Xmod$. Under this functor, every concept in the
sheaf-theoretic account has an exact counterpart in the variational account, and vice
versa.

\subsection{The Sheaf-Theoretic Account}

Let $\mathcal{M}$ be a topological space equipped with an open cover
$\mathcal{U} = \{U_i\}_{i \in \mathcal{I}}$. A presheaf $\mathcal{S}$ over $\mathcal{U}$
assigns to each $U_i$ a category $\mathcal{S}(U_i)$ of local sections, and to each
inclusion $U_j \subseteq U_i$ a restriction functor
$\rho_{ij} : \mathcal{S}(U_i) \to \mathcal{S}(U_j)$.

\begin{definition}[Gluing Condition]
A family of local sections $\{s_i \in \mathcal{S}(U_i)\}_{i \in \mathcal{I}}$ satisfies
the \emph{gluing condition} if for all $i, j \in \mathcal{I}$:
\[
\rho_{U_i, U_i \cap U_j}(s_i) = \rho_{U_j, U_i \cap U_j}(s_j) \quad \text{on }
U_i \cap U_j.
\]
A presheaf in which every compatible family of local sections admits a unique global
section $s \in \mathcal{S}(\mathcal{M})$ with $s|_{U_i} = s_i$ is called a \emph{sheaf}.
\end{definition}

\begin{definition}[Semantic Tear and Repair]
A \emph{semantic tear} at overlap $U_i \cap U_j$ is a failure of the gluing condition.
The discrepancy defines a \v{C}ech 1-cochain
$c_{ij} = \rho_{ij}(s_j) - \rho_{ji}(s_i)$. A \emph{repair} is a morphism in
$\mathcal{S}(U_i \cap U_j)$ that trivializes $c_{ij}$. The space of unresolved tears is
measured by $H^1(\mathcal{U}, \mathcal{S})$.
\end{definition}

\begin{definition}[Sheaf Repair Operator]
The \emph{sheaf repair operator} $\mathcal{R}_{\mathrm{sheaf}}$ maps a family of
inconsistent local sections to a minimally modified family satisfying the gluing condition,
where minimality is measured by a seam cost functional $\mathcal{L}_{\mathrm{seam}}$:
\[
\mathcal{R}_{\mathrm{sheaf}}(\{s_i\}) = \arg\min_{\{s'_i\}} \sum_{i < j}
\mathcal{L}_{\mathrm{seam}}(s'_i|_{U_{ij}}, s'_j|_{U_{ij}}) +
\sum_i \mathcal{L}_{\mathrm{local}}(s'_i, s_i).
\]
Global section existence corresponds to $\mathcal{L}_{\mathrm{seam}} = 0$ at the optimum.
\end{definition}

\subsection{The Variational Account}

The variational account centers on the consistency operator:
\[
\mathcal{C}([X]) = \arg\min_{[X'] \in \mathcal{A}}
\left[ \sum_{i \in \mathcal{I}} \|\Pi_i(X') - y_i\|_{\mathcal{Y}_i}^2
+ \mathcal{R}(X') \right]
\]
where $\Pi_i : \mathcal{X} \to \mathcal{Y}_i$ are projection operators, $y_i$ are
observations, and $\mathcal{R}$ is the RSVP-consistent regularizer encoding admissibility.
A fixed point $[X^*]$ satisfies $\Pi_i(X^*) \approx y_i$ for all $i$ within tolerance
$\epsilon_i$.

\subsection{The Dictionary}

The following correspondence identifies every object in the sheaf-theoretic account with
its exact counterpart in the variational account.

\begin{center}
\begin{tabular}{ll}
\toprule
\textbf{Sheaf-theoretic concept} & \textbf{Variational counterpart} \\
\midrule
Open set $U_i$ in cover $\mathcal{U}$ & Sensor modality $i$ with space $\mathcal{Y}_i$ \\
Local section $s_i \in \mathcal{S}(U_i)$ & Observation $y_i \in \mathcal{Y}_i$ \\
Restriction map $\rho_{ij}$ & Projection operator $\Pi_i$ \\
Gluing condition & Projection consistency $\Pi_i(X) = y_i$ \\
\v{C}ech 1-cochain $c_{ij}$ & Projection residual $\|\Pi_i(X) - y_i\|$ \\
Coboundary (trivialized cocycle) & Residual driven to zero by $\mathcal{C}$ \\
Global section $s \in \mathcal{S}(\mathcal{M})$ & Fixed point $[X^*] \in \Xmod$ \\
Cohomological obstruction $H^1 \ne 0$ & Projection degeneracy \\
Flat direction in $\mathcal{L}_{\mathrm{local}}$ & Regularizer flatness \\
Repair morphism trivializing $c_{ij}$ & One step of the refinement loop \\
Sheaf repair operator $\mathcal{R}_{\mathrm{sheaf}}$ & Consistency operator $\mathcal{C}$ \\
Stigmergic threshold & Contractivity condition on $\mathcal{C}$ \\
Self-maintaining growth & Convergence of refinement loop to $[X^*]$ \\
\bottomrule
\end{tabular}
\end{center}

\subsection{The Sheaf-Variational Equivalence}

\begin{definition}[Category of Projection Families]
Let $\mathbf{Proj}(\mathcal{X})$ be the category whose objects are pairs
$(\mathcal{I}, \{\Pi_i\}_{i \in \mathcal{I}})$ and whose morphisms are inclusions of
index sets compatible with the projection structure.
\end{definition}

\begin{definition}[Category of Covering Sheaves]
Let $\mathbf{Sh}(\mathcal{M})$ be the category whose objects are sheaves $\mathcal{S}$
over covers $\mathcal{U}$ of $\mathcal{M}$, and whose morphisms are morphisms of sheaves
over refinements of covers.
\end{definition}

\begin{theorem}[Sheaf-Variational Equivalence]
\label{thm:equivalence}
There exists a functor
\[
\Psi : \mathbf{Proj}(\mathcal{X}) \longrightarrow \mathbf{Sh}(\Xmod)
\]
such that:
\begin{enumerate}
\item[(i)] $\Psi$ maps each projection family $\{\Pi_i\}_{i \in \mathcal{I}}$ to a sheaf
$\mathcal{S}_\Pi$ over a cover of $\Xmod$ whose local sections are observation-consistent
field configurations.
\item[(ii)] The gluing condition of $\mathcal{S}_\Pi$ is equivalent to the projection
consistency condition $\Pi_i(X) = y_i$ for all $i$.
\item[(iii)] The consistency operator $\mathcal{C}$ is the variational realization of the
sheaf repair operator $\mathcal{R}_{\mathrm{sheaf}}$ under $\Psi$.
\item[(iv)] The fixed point $[X^*]$ of $\mathcal{C}$ corresponds under $\Psi$ to the
unique global section of $\mathcal{S}_\Pi$.
\item[(v)] The cohomological obstruction $H^1(\mathcal{U}, \mathcal{S}_\Pi) = 0$ if and
only if the projection family is identifying on $\mathcal{F}$.
\end{enumerate}
\end{theorem}

\begin{proof}
\textbf{Construction of $\Psi$.} Given a projection family
$(\mathcal{I}, \{\Pi_i\})$, construct the cover
$\mathcal{U} = \{U_i\}_{i \in \mathcal{I}}$ of $\Xmod$ where
\[
U_i = \{ [X] \in \mathcal{A} \mid \|\Pi_i(X) - y_i\|_{\mathcal{Y}_i} \le \epsilon_i \}.
\]
Define the sheaf $\mathcal{S}_\Pi$ by assigning to each $U_i$:
\[
\mathcal{S}_\Pi(U_i) = U_i \cap \mathcal{A}.
\]
For inclusions $U_j \subseteq U_i$, the restriction map $\rho_{ij}$ is the inclusion
$U_j \cap \mathcal{A} \hookrightarrow U_i \cap \mathcal{A}$.

\textbf{Part (i).} By construction, $\mathcal{S}_\Pi(U_i)$ consists of admissible
configurations whose $i$-th projection is within tolerance of $y_i$. $\checkmark$

\textbf{Part (ii).} A family $\{[X_i] \in \mathcal{S}_\Pi(U_i)\}$ satisfies the gluing
condition iff $[X_i] = [X_j]$ on $U_i \cap U_j$ for all $i,j$, which is precisely the
condition that a single $[X] \in \mathcal{A}$ satisfies $\Pi_i(X) = y_i$ for all $i$
simultaneously. $\checkmark$

\textbf{Part (iii).} The objective of $\mathcal{R}_{\mathrm{sheaf}}$ translates as:
\[
\sum_{i<j} \mathcal{L}_{\mathrm{seam}} + \mathcal{L}_{\mathrm{local}}
\;\longleftrightarrow\;
\sum_i \|\Pi_i(X) - y_i\|^2 + \mathcal{R}(X),
\]
which is exactly the objective of $\mathcal{C}$. $\checkmark$

\textbf{Part (iv).} A fixed point $[X^*]$ of $\mathcal{C}$ satisfies
$[X^*] \in U_i$ for all $i$, giving a local section $s_i = [X^*]|_{U_i}$ that trivially
satisfies the gluing condition. Uniqueness follows from Theorem~\ref{thm:identifiability}.
$\checkmark$

\textbf{Part (v).} $H^1(\mathcal{U}, \mathcal{S}_\Pi)$ measures failure of compatible
local sections to glue into a global section, which is exactly projection degeneracy by
part (ii). $\checkmark$

\textbf{Functoriality.} Adding sensors corresponds to refining the cover; the refinement
map defines a natural transformation establishing functoriality. $\square$
\end{proof}

\subsection{Consequences of the Equivalence}

\textbf{Consequence 1: The two failure modes are cohomological.} Projection degeneracy is
$H^1 \ne 0$; regularizer flatness is degenerate $H^0$. Adding sensors can kill $H^1$
classes; strengthening the regularizer thins the stalks and kills degenerate $H^0$.

\textbf{Consequence 2: Dynamical disambiguation is cohomology in time.} The dynamical
admissibility set $\Gamma$ extends the cover into the temporal direction. Trajectory
separation kills residual $H^1$ classes that spatial sensors cannot. The full feasible set
$\mathcal{F} = \Omega_{\mathrm{obs}} \cap \Gamma$ is the global section of a sheaf over a
spatio-temporal cover.

\textbf{Consequence 3: Forward and inverse are dual.} The forward direction studies the
repair dynamics of $\mathcal{S}_\Pi$ from inside---how sections evolve under stigmergic
reinforcement and homeorhetic flow. The inverse direction studies the existence and
uniqueness of the global section from outside. Both study the same sheaf.

\subsection{The Entropy Identification}

The RSVP entropy field $S$ appears in two roles throughout this framework: as the
entropy component of the field triple $(\Phi, \mathbf{v}, S)$, and as a measure of
unresolved obstruction. Under the Sheaf-Variational Equivalence these become one.

The observation-consistent set $\Omega_{\mathrm{obs}}$ carries a natural measure induced
by the Fisher information metric of the noise models $\{\mathcal{N}_i\}$ (developed in
Appendix~\ref{app:fisher}). The volume of $\Omega_{\mathrm{obs}}$ under this metric
measures residual uncertainty. Under the equivalence:
\[
\log \vol_{\mathrm{Fisher}}(\Omega_{\mathrm{obs}})
\;\longleftrightarrow\;
\log \dim H^0(\mathcal{U}, \mathcal{S}_\Pi)
\;\longleftrightarrow\;
S_{\mathrm{RSVP}}(X^*).
\]
All three measure residual degrees of freedom after all constraints have been applied.
Closure corresponds to each collapsing simultaneously. The RSVP entropy field is not
merely a bookkeeping device---it is the local density of cohomological obstruction, and
its minimization under $\mathcal{C}$ is the physical expression of sheaf repair.

%% ============================================================
\section{The World-State Manifold and Its Symmetry Group}
\label{sec:worldstate}
%% ============================================================

\preliminary{
\textbf{This section develops the minimal mathematical infrastructure that Appendix B
requires: the configuration space $\mathcal{X}$, its symmetry group $G$, the quotient
$\mathcal{X}/{\sim}$, and the admissible set $\mathcal{A}$.}

\medskip
\textbf{3.1 The Configuration Space.}
Defines $\mathcal{X} = H^1(\Omega) \times H^1(\Omega;\mathbb{R}^3) \times L^2(\Omega)$
over compact domain $\Omega \subset \mathbb{R}^3$. Specifies the product Sobolev norm.
States that $\mathcal{X}$ is a separable Hilbert space and hence also a Fréchet manifold
(this is elaborated in Appendix~\ref{app:frechet}). Identifies the three field components
with physical quantities: $\Phi$ as ecosystem scalar density and civilizational legitimacy
potential, $\mathbf{v}$ as resource flow and infrastructure transport vectors, $S$ as
entropy production and complexity budget. Gives explicit examples of each component in
ecological and civilizational contexts to ground the abstraction before the symmetry group
is introduced.

\medskip
\textbf{3.2 The Symmetry Group.}
Defines $G$ as the group acting on $\mathcal{X}$ comprising: (a) gauge redundancies of
the field representation---rephasing and rescaling transformations that leave physical
observables invariant; (b) smooth coordinate reparametrizations of $\Omega$---diffeomorphisms
that relabel spatial positions without changing field content; (c) observationally
equivalent reconfigurations of the civilizational layer---distinct infrastructure or
demographic configurations that produce identical sensor outputs. States that the action
of $G$ is smooth, proper, and isometric with respect to the product Sobolev metric, so
that $\mathcal{X}/G$ is Hausdorff and inherits a natural metric. Distinguishes $G$-equivalence
from observational equivalence $\sim$: the former is intrinsic to the field representation,
the latter is extrinsic and depends on the projection family. Notes that $\sim$ is coarser
than $G$-equivalence in general and will collapse to $G$-equivalence when the projection
family is complete (as established in Section~\ref{sec:projections}).

\medskip
\textbf{3.3 The Admissible Set.}
Defines $\mathcal{A} \subset \mathcal{X}/{\sim}$ as the set of physically admissible
equivalence classes: those satisfying the RSVP field equations, ecological plausibility
constraints, and civilizational feasibility constraints. States Assumption A4 (geodesic
convexity and strict convexity of $\mathcal{R}$ on $\mathcal{A}$) and gives the physical
motivation for each condition: geodesic convexity ensures that admissible interpolations
respect symmetry, lower semicontinuity ensures admissibility is stable under limits,
coercivity ensures physically implausible configurations are excluded rather than merely
penalized. Identifies $\mathcal{A}$ as a closed convex subset of the Banach realization
$\mathcal{B}$ defined in Assumption B1 of Appendix~\ref{app:identifiability}.

\medskip
\textbf{3.4 The RSVP Field Equations.}
States the plenum field equations governing time evolution of $(\Phi, \mathbf{v}, S)$:
scalar continuity, vector momentum, and entropy production equations. These are the
equations that $\mathcal{R}_{\mathrm{RSVP}}$ penalizes violation of. Notes that solutions
to these equations form the dynamical admissibility set $\Gamma$ in the zero-noise limit
(Section~\ref{sec:disambiguation}). Does not derive the equations---they are taken as
given from the RSVP framework---but states them precisely enough that $\mathcal{R}_{\mathrm{RSVP}}$
can be defined as a specific $L^2$ norm of residuals.

\medskip
\textbf{3.5 The Regularizer.}
Defines
\[
\mathcal{R}(X) = \mathcal{R}_{\mathrm{RSVP}}(X)
+ \lambda_{\mathrm{eco}}\,\mathcal{R}_{\mathrm{eco}}(X)
+ \lambda_{\mathrm{civ}}\,\mathcal{R}_{\mathrm{civ}}(X)
\]
with each component defined explicitly. $\mathcal{R}_{\mathrm{RSVP}}$: squared $L^2$
residuals of the three RSVP field equations. $\mathcal{R}_{\mathrm{eco}}$: zero on
ecologically admissible configurations (satisfying biomass conservation, succession
ordering, watershed feasibility), strictly positive off it. $\mathcal{R}_{\mathrm{civ}}$:
zero on civilizationally admissible configurations (satisfying demographic feasibility,
infrastructure stress bounds, resource flow conservation), strictly positive off it.
Emphasizes that the decomposition is notational: $\mathcal{R}$ is a single functional
encoding coupled closure, not a product of independent penalties.
}

%% ============================================================
\section{Sensor Modalities as Projection Operators}
\label{sec:projections}
%% ============================================================

\subsection{Purpose of This Section}

Section~\ref{sec:equivalence} established that a projection family $\{\Pi_i\}_{i \in \mathcal{I}}$
and a sheaf $\mathcal{S}_\Pi$ over $\Xmod$ are the same structure viewed from outside and
inside respectively. What remains is to specify what the projection operators actually
are---not as abstract bounded linear maps, but as concrete operators whose domains,
codomains, and kernels are determined by the physics of each sensor modality.

This section defines each projection operator precisely enough that Assumption A1 can be
verified directly. It establishes the RSVP correspondence: the identification of each
modality with a specific component or functional of the field triple $(\Phi, \mathbf{v}, S)$.
And it proves the completeness proposition: the condition under which the projection family
collectively spans the full RSVP triple.

Nothing in this section is motivated by sensor engineering. The operators are defined as
they are because those definitions are what the theorems in Sections~\ref{sec:identifiability}
and \ref{sec:disambiguation} require.

\subsection{The Configuration Space and Its Topology}

Let $\Omega \subset \mathbb{R}^3$ be a compact domain with smooth boundary. The
configuration space is
\[
\mathcal{X} = H^1(\Omega) \times H^1(\Omega; \mathbb{R}^3) \times L^2(\Omega)
\]
with product Sobolev norm $\|X\|_{\mathcal{X}}^2 = \|\Phi\|_{H^1}^2 + \|\mathbf{v}\|_{H^1}^2 + \|S\|_{L^2}^2$.
All projection operators defined below are required to descend to the quotient $\Xmod$:
they must be constant on $G$-orbits.

\subsection{Infrared Projection}

\begin{definition}[Infrared Projection]
The infrared projection $\Pi_{\mathrm{IR}} : \mathcal{X} \to \mathcal{Y}_{\mathrm{IR}}$
is defined by
\[
\Pi_{\mathrm{IR}}(X) = \left( \Phi|_{\partial\Omega},\;
\nabla\Phi \cdot \hat{n}|_{\partial\Omega} \right)
\]
where $\hat{n}$ is the outward unit normal to $\partial\Omega$ and
$\mathcal{Y}_{\mathrm{IR}} = L^2(\partial\Omega) \times L^2(\partial\Omega)$.
\end{definition}

\begin{lemma}
$\Pi_{\mathrm{IR}}$ is a bounded linear operator from $\mathcal{X}$ to $\mathcal{Y}_{\mathrm{IR}}$.
\end{lemma}
\begin{proof}
The trace operator $\Phi \mapsto \Phi|_{\partial\Omega}$ is bounded from $H^1(\Omega)$
to $L^2(\partial\Omega)$ by the Sobolev trace theorem. The normal derivative is bounded
from $H^1(\Omega)$ to $H^{-1/2}(\partial\Omega) \hookrightarrow L^2(\partial\Omega)$ for
smooth $\partial\Omega$. $\square$
\end{proof}

\textbf{RSVP correspondence.} $\Pi_{\mathrm{IR}}$ constrains boundary values and outward
flux of $\Phi$: energy density at the surface and rate of energy exchange. It does not
directly constrain $\mathbf{v}$ or $S$ in the interior.

\subsection{LiDAR Projection}

\begin{definition}[LiDAR Projection]
Fix a threshold $\phi_0 > 0$. The LiDAR projection
$\Pi_{\mathrm{LiDAR}} : \mathcal{X} \to \mathcal{Y}_{\mathrm{LiDAR}}$ is defined by
\[
\Pi_{\mathrm{LiDAR}}(X) = \mathbf{1}_{\{\Phi \ge \phi_0\}}
\quad \text{together with} \quad
\Pi_{\mathrm{LiDAR}}^{\perp}(X) =
\frac{\nabla\Phi}{|\nabla\Phi|}\bigg|_{\{\Phi = \phi_0\}}.
\]
The observation space is
$\mathcal{Y}_{\mathrm{LiDAR}} = L^2(\Omega;\{0,1\}) \times L^2(\{\Phi=\phi_0\};S^2)$.
\end{definition}

\begin{lemma}
$\Pi_{\mathrm{LiDAR}}$ is bounded for $\phi_0$ a regular value of $\Phi$.
\end{lemma}
\begin{proof}
For $\phi_0$ a regular value, $\{\Phi = \phi_0\}$ is smooth by the implicit function
theorem in $H^1$. The indicator is bounded in $L^2(\Omega)$ by domain volume. The normal
field is bounded by the Sobolev embedding $H^1 \hookrightarrow C^0$ in three dimensions.
$\square$
\end{proof}

\textbf{RSVP correspondence.} $\Pi_{\mathrm{LiDAR}}$ constrains the geometry of level
sets of $\Phi$ and normal directions to those surfaces, which simultaneously constrain
boundary conditions on $\mathbf{v}$: flows cannot cross material boundaries.

\subsection{RGB/CCD Projection}

\begin{definition}[RGB Projection]
Let $\mathbf{c} \in S^2$ be the camera viewing direction and $f : \mathbb{R} \to \mathbb{R}^3$
a material reflectance function. The RGB projection
$\Pi_{\mathrm{RGB}} : \mathcal{X} \to \mathcal{Y}_{\mathrm{RGB}}$ is defined by
\[
\Pi_{\mathrm{RGB}}(X) = \int_0^\infty f(\Phi(\mathbf{x}(t)))\,
e^{-\int_0^t \Phi(\mathbf{x}(s)) ds}\, dt
\]
where $\mathbf{x}(t) = \mathbf{x}_0 + t\mathbf{c}$ and
$\mathcal{Y}_{\mathrm{RGB}} = L^2(\mathcal{V};\mathbb{R}^3)$ over camera pixels $\mathcal{V}$.
\end{definition}

\begin{lemma}
$\Pi_{\mathrm{RGB}}$ is bounded when $f$ is Lipschitz and $\Phi \ge 0$.
\end{lemma}
\begin{proof}
The exponential attenuation is bounded by 1. The integrand is bounded by $\|f\|_\infty$,
which is finite by Lipschitz continuity and $H^1 \hookrightarrow L^\infty$ in three
dimensions. $\square$
\end{proof}

\textbf{RSVP correspondence.} $\Pi_{\mathrm{RGB}}$ constrains $\Phi$ along rays---a
projective complement to $\Pi_{\mathrm{LiDAR}}$'s level-set constraint. Their kernels
overlap partially but not completely.

\subsection{Audio Projection}

\begin{definition}[Audio Projection]
Let $\mathbf{x}_m \in \Omega$ for $m = 1,\ldots,M$ be microphone positions. The audio
projection $\Pi_{\mathrm{audio}} : \mathcal{X} \times [0,T] \to \mathcal{Y}_{\mathrm{audio}}$
is
\[
\Pi_{\mathrm{audio}}(X)(t) =
\left( \partial_t S(\mathbf{x}_m, t) + \nabla \cdot \mathbf{v}(\mathbf{x}_m, t)
\right)_{m=1}^M
\]
with $\mathcal{Y}_{\mathrm{audio}} = L^2([0,T];\mathbb{R}^M)$.
\end{definition}

\textbf{RSVP correspondence.} $\Pi_{\mathrm{audio}}$ is the only modality that directly
constrains $\partial_t S$---the rate of entropy production---and $\nabla \cdot \mathbf{v}$,
encoding sources and sinks of flow. Hidden processes invisible to all spatial sensors are
audible as entropy production events.

\subsection{Photonic Strip Projection}

\begin{definition}[Photonic Strip Projection]
Let $\gamma : [0,L] \to \Omega$ parametrize the strip by arc length. The photonic strip
projection $\Pi_{\mathrm{strip}} : \mathcal{X} \times [0,T] \to \mathcal{Y}_{\mathrm{strip}}$
is
\[
\Pi_{\mathrm{strip}}(X)(s,t) =
\left( \Phi(\gamma(s),t),\;
\mathbf{v}(\gamma(s),t) \cdot \gamma'(s),\;
\partial_t \Phi(\gamma(s),t) \right)
\]
with $\mathcal{Y}_{\mathrm{strip}} = L^2([0,L] \times [0,T];\mathbb{R}^3)$.
\end{definition}

\textbf{RSVP correspondence.} The photonic strip provides continuous temporal records of
$\Phi$, tangential $\mathbf{v}$, and $\partial_t \Phi$ along a fixed transect: the only
modality directly constraining temporal evolution of $\Phi$ at interior points.

\subsection{The Completeness Proposition}

\begin{definition}[Spanning Condition]
A projection family $\{\Pi_i\}_{i \in \mathcal{I}}$ is \emph{field-spanning} if
\[
\ker \Ftil = \{ [X] \in \Xmod \mid \Pi_i(X) = 0 \;\forall i \in \mathcal{I} \} = \{[0]\}.
\]
\end{definition}

\begin{proposition}[Completeness of the Five-Modality Family]
\label{prop:completeness}
The projection family $\{\Pi_{\mathrm{IR}}, \Pi_{\mathrm{LiDAR}}, \Pi_{\mathrm{RGB}},
\Pi_{\mathrm{audio}}, \Pi_{\mathrm{strip}}\}$ is field-spanning on $\Xmod$.
\end{proposition}

\begin{proof}
Suppose $[X] \in \ker \Ftil$. $\Pi_{\mathrm{IR}}(X) = 0$ gives
$\Phi|_{\partial\Omega} = 0$. $\Pi_{\mathrm{LiDAR}}(X) = 0$ gives $\Phi < \phi_0$
a.e. $\Pi_{\mathrm{strip}}(X) = 0$ along strips dense in $\Omega$ gives $\Phi = 0$
on a dense subset; by $H^1$ continuity, $\Phi = 0$ on $\Omega$. With $\Phi = 0$,
strip vanishing also gives $\mathbf{v} \cdot \gamma' = 0$ along all strips; with three
non-parallel strip directions, $\mathbf{v} = 0$. Then $\Pi_{\mathrm{audio}}(X) = 0$
with $\mathbf{v} = 0$ gives $\partial_t S = 0$ at dense microphone positions, so $S$
is time-constant; RSVP entropy constraints force $S = 0$. Hence $[X] = [0]$. $\square$
\end{proof}

\subsection{Degrees of Freedom Table}

\begin{center}
\begin{tabular}{llll}
\toprule
\textbf{Modality} & \textbf{Field constrained} & \textbf{Where} &
\textbf{New DOF closed} \\
\midrule
Infrared & $\Phi$, $\nabla\Phi \cdot \hat{n}$ & $\partial\Omega$ &
Boundary energy, flux \\
LiDAR & Level sets of $\Phi$, normals & Interior surfaces &
Geometry; $\mathbf{v}$ BCs \\
RGB & $\Phi$ along rays & Lines of sight &
Interior density \\
Audio & $\partial_t S$, $\nabla \cdot \mathbf{v}$ & Point measurements &
Temporal entropy, flow sources \\
Photonic strip & $\Phi$, $\mathbf{v} \cdot \gamma'$, $\partial_t \Phi$ &
Transects over time & Temporal $\Phi$; tangential flow \\
\bottomrule
\end{tabular}
\end{center}

\subsection{Each Modality as a Chart}

Under the Sheaf-Variational Equivalence, each $\Pi_i$ defines a chart
$U_i = \{[X] \in \mathcal{A} \mid \|\Pi_i(X) - y_i\| \le \epsilon_i\}$ on $\Xmod$.
The five modalities define five charts. Completeness ensures the cover is fine enough
to support a non-degenerate sheaf. Sensor design is sheaf design: choosing modalities
is choosing the cover of $\Xmod$.

%% ============================================================
\section{The Feasible Set and Its Geometry}
\label{sec:feasible}
%% ============================================================

\preliminary{
\textbf{This section defines and analyzes the geometric object whose dimensional collapse
is the operative notion of closure throughout the paper.}

\medskip
\textbf{4.1 The Observation-Consistent Set.}
Defines $\Omega_{\mathrm{obs}} = \{[X] \in \mathcal{A} \mid \|\Pi_i(X) - y_i\|_{\mathcal{Y}_i}
\le \epsilon_i\; \forall i\}$ and proves it is non-empty (Lemma 5.1), closed, and convex
(Lemma 5.2), using results already established in Section~\ref{sec:identifiability}.
Gives geometric intuition: $\Omega_{\mathrm{obs}}$ is the intersection of five ``tubes''
in $\Xmod$, one per sensor modality, each of radius $\epsilon_i$ around the preimage
of $y_i$ under $\Pi_i$.

\medskip
\textbf{4.2 The Dynamical Admissibility Set.}
Introduces $\Gamma$ as a named object whose full definition appears in
Section~\ref{sec:disambiguation}. States its key property here: $\Gamma$ is a subset of
$\mathcal{A}$ consisting of configurations that lie on admissible trajectories of the world
dynamics. Notes that $\Gamma \subseteq \mathcal{A}$ by forward-invariance (Lemma 6.1), and
that in the deterministic limit $\Gamma^0$ is the union of integral curves of the vector
field $F$.

\medskip
\textbf{4.3 The Feasible Set.}
Defines $\mathcal{F} = \Omega_{\mathrm{obs}} \cap \Gamma$ as the primary geometric object
of the paper. Notes that $\mathcal{F}$ is the correct domain for both theorems: the
Identifiability Theorem requires $\Ftil$ to be injective on $\mathcal{F}$, not on all of
$\Xmod$; the Disambiguation Theorem requires trajectory separation on $\mathcal{F}$, not
on $\Omega_{\mathrm{obs}}$ alone.

\medskip
\textbf{4.4 Topology of the Feasible Set.}
Analyzes the topology of $\mathcal{F}$ as a function of the projection family and the
$\epsilon_i$ profile. The key quantity is $\dim \mathcal{F}$: closure corresponds to
$\dim \mathcal{F} \to 0$. States and proves the Feasible Set Contraction theorem: under
monotone addition of sensors with independent noise models, $\dim \mathcal{F}$ is
non-increasing, with strict decrease when the new sensor constrains at least one degree of
freedom not already closed by $\{\Pi_1,\ldots,\Pi_{i-1}\}$. Connects this to the
completeness result of Section~\ref{sec:projections}: the five-modality family is complete
iff it can in principle drive $\dim \mathcal{F}$ to zero.

\medskip
\textbf{4.5 Information Geometry of the Feasible Set.}
Introduces the Fisher information metric on $\mathcal{F}$ induced by the noise models
$\{\mathcal{N}_i\}$. Under this metric, $\mathcal{F}$ is a Riemannian manifold and its
volume $\vol_{\mathrm{Fisher}}(\mathcal{F})$ is the measure of residual uncertainty.
States that closure corresponds to $\vol_{\mathrm{Fisher}}(\mathcal{F}) \to 0$. Identifies
this volume with the RSVP entropy and cohomological obstruction as established in
Section~\ref{sec:equivalence}. Full development of the Fisher geometry is deferred to
Appendix~\ref{app:fisher}.

\medskip
\textbf{4.6 The Sheaf Interpretation.}
Under the Sheaf-Variational Equivalence, $\mathcal{F}$ is the intersection of stalks
across all charts of the cover $\mathcal{U}$. Closure $|\mathcal{F}| = 1$ is exactly the
sheaf condition: there exists a unique global section consistent across all charts.
Non-existence of a global section is not a failure mode of the system---it is a
falsification signal indicating that the model class $\mathcal{A}$ needs revision. This
is an instance of the general principle from sheaf theory: non-trivial $H^1$ is
informative, not merely obstructive.
}

%% ============================================================
\section{The Consistency Operator and the Identifiability Theorem}
\label{sec:identifiability}
%% ============================================================

\subsection{Setup and Standing Assumptions}

Let $\mathcal{X} = H^1(\Omega) \times H^1(\Omega;\mathbb{R}^3) \times L^2(\Omega)$
with the product Sobolev topology (Appendix~\ref{app:frechet}). A symmetry group $G$
acts smoothly, properly, and isometrically on $\mathcal{X}$; the quotient $\Xmod$ is
Hausdorff. We work in $\Xmod$ throughout.

\begin{assumption}[Projection Regularity]
\label{ass:regularity}
Each $\Pi_i : \mathcal{X} \to \mathcal{Y}_i$ is a bounded linear map to a separable
Hilbert space $\mathcal{Y}_i$.
\end{assumption}

\begin{assumption}[Noise Model]
\label{ass:noise}
Each observation satisfies $y_i = \Pi_i(X^*) + \eta_i$ where $\|\eta_i\|_{\mathcal{Y}_i}
\le \epsilon_i$ almost surely.
\end{assumption}

\begin{assumption}[Regularizer Structure]
\label{ass:regularizer}
The regularizer takes the form
\[
\mathcal{R}(X) = \mathcal{R}_{\mathrm{RSVP}}(X)
+ \lambda_{\mathrm{eco}}\,\mathcal{R}_{\mathrm{eco}}(X)
+ \lambda_{\mathrm{civ}}\,\mathcal{R}_{\mathrm{civ}}(X)
\]
where $\mathcal{R}_{\mathrm{RSVP}}$ measures violation of the RSVP field equations in
$L^2$, $\mathcal{R}_{\mathrm{eco}}$ is zero on ecologically admissible configurations
and strictly positive off them, and $\mathcal{R}_{\mathrm{civ}}$ is zero on
civilizationally admissible configurations and strictly positive off them.
\end{assumption}

\begin{assumption}[Admissibility Region]
\label{ass:admissibility}
The admissible set $\mathcal{A} = \{[X] \in \Xmod \mid \mathcal{R}(X) < \infty\}$
is a closed convex subset of $\Xmod$, and $\mathcal{R}$ restricted to $\mathcal{A}$
is strictly convex along geodesics.
\end{assumption}

\begin{assumption}[Banach Realization]
\label{ass:banach}
There exists a Banach space $\mathcal{B}$ and continuous embedding
$\iota : \mathcal{X} \hookrightarrow \mathcal{B}$ such that $\mathcal{A}$ is closed and
convex in $\mathcal{B}$, each $\Pi_i$ extends to a bounded linear operator on
$\mathcal{B}$, and $\mathcal{R}$ extends to a proper lower semicontinuous functional on
$\mathcal{B}$.
\end{assumption}

In practice one takes $\mathcal{B} = \mathcal{X}$ itself, which is already a Hilbert
space. The Banach realization is stated separately to make the compactness and weak
convergence arguments in Appendix~\ref{app:identifiability} self-contained.

\subsection{The Feasible Set}

\begin{definition}[Observation-Consistent Set]
\[
\Omega_{\mathrm{obs}} = \{ [X] \in \mathcal{A} \mid
\|\Pi_i(X) - y_i\|_{\mathcal{Y}_i} \le \epsilon_i \;\; \forall i \in \mathcal{I} \}.
\]
\end{definition}

\begin{lemma}[Non-emptiness]
Under Assumptions~\ref{ass:regularity} and \ref{ass:noise}, $X^* \in \Omega_{\mathrm{obs}}$
almost surely.
\end{lemma}
\begin{proof}
$\|y_i - \Pi_i(X^*)\| = \|\eta_i\| \le \epsilon_i$ almost surely by Assumption~\ref{ass:noise}.
$\square$
\end{proof}

\begin{lemma}[Closedness and Convexity]
Under Assumptions~\ref{ass:regularity} and \ref{ass:admissibility},
$\Omega_{\mathrm{obs}}$ is a closed convex subset of $\mathcal{A}$.
\end{lemma}
\begin{proof}
Each constraint set is closed (continuity of $\Pi_i$) and convex (linearity of $\Pi_i$,
convexity of norm ball). Intersection of closed convex sets is closed and convex.
$\square$
\end{proof}

The dynamically admissible set $\Gamma$ is introduced as a named object here; its full
definition and construction are given in Section~\ref{sec:disambiguation}. The
\emph{feasible set} is $\mathcal{F} = \Omega_{\mathrm{obs}} \cap \Gamma$.

\subsection{The Consistency Operator}

\begin{definition}[Consistency Operator]
\[
\mathcal{C}(X) = \arg\min_{[X'] \in \mathcal{A}}
\left[ \sum_{i \in \mathcal{I}} \|\Pi_i(X') - y_i\|_{\mathcal{Y}_i}^2
+ \mathcal{R}(X') \right].
\]
\end{definition}

$\mathcal{C}$ maps $\mathcal{A}$ to $\mathcal{A}$: the optimization is constrained to the
admissible region. A configuration minimizing projection error but violating RSVP equations
or ecological plausibility is not a candidate world-state. The regularizer defines the
domain of the operator, not a preference over an unconstrained space.

\begin{lemma}[Well-definedness]
Under Assumptions~\ref{ass:regularity}--\ref{ass:banach}, $\mathcal{C}(X)$ exists and is
unique for every $[X] \in \mathcal{A}$.
\end{lemma}
\begin{proof}
The objective $f(X') = \sum_i \|\Pi_i(X') - y_i\|^2 + \mathcal{R}(X')$ is continuous,
strictly convex (sum of convex functions plus strictly convex $\mathcal{R}$), and coercive
on $\mathcal{A}$. A continuous strictly convex coercive function on a closed convex set
attains its minimum at a unique point. Full proof in Appendix~\ref{app:identifiability}.
$\square$
\end{proof}

\subsection{The Identifiability Theorem}

The induced map of the projection family on the quotient is:
\[
\Ftil : \Xmod \longrightarrow \prod_{i \in \mathcal{I}} \mathcal{Y}_i, \qquad
[X] \mapsto (\Pi_i(X))_{i \in \mathcal{I}}.
\]
By definition of $\sim$, $\Ftil$ is well-defined.

\begin{theorem}[Yarncrawler Identifiability]
\label{thm:identifiability}
Under Assumptions~\ref{ass:regularity}--\ref{ass:banach}, the consistency operator
$\mathcal{C}$ has a unique fixed point $[X^*] \in \mathcal{F}$ if and only if $\Ftil$
is injective on $\mathcal{F}$.
\end{theorem}

\begin{proof}
$(\Rightarrow)$ Suppose $\mathcal{C}$ has a unique fixed point $[X^*] \in \mathcal{F}$.
If $\Ftil$ were not injective on $\mathcal{F}$, there exist distinct $[X],[X'] \in \mathcal{F}$
with $\Pi_i(X) = \Pi_i(X')$ for all $i$. Both achieve zero projection residual in the
objective. The objective reduces to $\mathcal{R}(X)$ vs.\ $\mathcal{R}(X')$. Either
$\mathcal{R}(X) = \mathcal{R}(X')$---violating strict convexity of $\mathcal{R}$ on
$\mathcal{A}$ (Assumption~\ref{ass:admissibility})---or one has strictly smaller
regularizer value, making the other not a minimizer. Either way contradicts uniqueness.

$(\Leftarrow)$ If $\Ftil$ is injective on $\mathcal{F}$, then $\Omega_{\mathrm{obs}} \cap
\mathcal{F}$ contains at most one point. By Lemma 5.1, $[X^*] \in \Omega_{\mathrm{obs}}$;
by definition $[X^*] \in \Gamma$; hence $[X^*] \in \mathcal{F}$. Injectivity forces
$|\mathcal{F}| = 1$. At $[X^*]$ all projection residuals are within tolerance and
$\mathcal{R}(X^*)$ is minimized over the singleton $\mathcal{F}$. Hence
$\mathcal{C}([X^*]) = [X^*]$. $\square$
\end{proof}

\subsection{The Corollary on Diagnosable Failure Modes}

\begin{corollary}
\label{cor:failures}
Identifiability fails in exactly two geometrically distinguishable ways:

\emph{(i) Projection degeneracy}: $\Ftil$ is not injective on $\mathcal{F}$. Two distinct
admissible world-states produce identical observations. This failure is intrinsic to the
sensor family and cannot be resolved by the regularizer.

\emph{(ii) Regularizer flatness}: $\mathcal{R}$ is not strictly convex on $\mathcal{F}$.
The regularizer assigns equal admissibility to configurations the sensors cannot
distinguish. This failure is intrinsic to the model class.

These failure modes are geometrically distinguishable: projection degeneracy corresponds
to a non-trivial kernel of $\Ftil|_{\mathcal{F}}$; regularizer flatness corresponds to a
flat direction of $\mathcal{R}$ within $\mathcal{F}$. They can coexist.
\end{corollary}

\subsection{Why the Quotient Structure Is Not Optional}

In a standard inverse problem one works in $\mathcal{X}$ directly. But in the Yarncrawler
setting, two configurations related by a gauge transformation or coordinate
reparametrization represent the same physical situation. Working in $\mathcal{X}$ rather
than $\Xmod$ would make uniqueness false by construction---there would always be a
continuum of configurations achieving the same objective, related by the action of $G$.

An $L^2$ uniqueness statement conflates field-theoretic gauge redundancy with measurement
noise. A measure-theoretic statement collapses the geometry of $\Xmod$ into observational
statistics and loses the regularizer's role as an admissibility condition. The
symmetry-group formulation is not a choice among equivalent formalisms; it is the only
formulation that preserves the semantic commitments of the framework.

%% ============================================================
\section{Trajectory Learning and the Dynamical Admissibility Set}
\label{sec:disambiguation}
%% ============================================================

\subsection{The Central Obstruction}

Section~\ref{sec:identifiability} established that $\mathcal{C}$ has a unique fixed point
iff $\Ftil$ is injective on $\mathcal{F}$. This condition is clean but generically unmet.
Under realistic sensing, $\Omega_{\mathrm{obs}}$ is high-dimensional. Projection degeneracy
is the default condition, not a pathological case.

The resolution cannot come from within the static framework. What is required is a
constraint of a different kind---one operating not on the configuration $X$ at a single
time but on the trajectory through which $X$ was reached. The dynamical admissibility set
$\Gamma$ is not an additional observation; it is a restriction on which configurations are
reachable under the dynamics of the world. Two configurations in $\Omega_{\mathrm{obs}}$
that are observationally indistinguishable at one moment may evolve under those dynamics
into configurations that are distinguishable at a later moment.

\subsection{The Path Space and the Stochastic Kernel}

\begin{definition}[Discrete Path Space]
$\mathcal{P}_\Delta = \{ \mathbf{X} = ([X_0],[X_1],\ldots,[X_N]) \mid [X_k] \in \mathcal{A},\;
N = \lfloor T/\Delta \rfloor \}$.
\end{definition}

\begin{definition}[Stochastic Transition Kernel]
$T_\Delta : \mathcal{A} \to \mathcal{P}(\mathcal{A})$, where for each $[X] \in \mathcal{A}$,
$T_\Delta([X],\cdot)$ specifies the distribution over one-step successors under the world
dynamics at resolution $\Delta$.
\end{definition}

\begin{definition}[Dynamically Admissible Set]
\[
\Gamma = \left\{ [X] \in \mathcal{A} \;\middle|\;
\exists\, [X_0] \in \mathcal{A},\; \exists\, k \ge 0 \text{ s.t. }
\mu_{X_0}^{T_\Delta}([X_k] = [X]) > 0 \right\}.
\]
$\Gamma$ is the support of the union of all admissible path measures. It is a geometric
object---a subset of $\mathcal{A}$---not a probabilistic one.
\end{definition}

\begin{lemma}[Invariance of $\Gamma$]
If $T_\Delta$ is time-homogeneous and $\mathcal{A}$ is forward-invariant under $T_\Delta$,
then $\Gamma \subseteq \mathcal{A}$ and $\Gamma$ is forward-invariant.
\end{lemma}

\subsection{The Deterministic Limit and Tangent Bundle Constraint}

\begin{assumption}[Generator Regularity]
\label{ass:generator}
There exists a Lipschitz vector field $F(\cdot\,;\theta) : \mathcal{A} \to T\mathcal{A}$
such that $T_\Delta([X],\cdot) \Rightarrow \delta_{[X]+\Delta F([X];\theta)}$ weakly as
$\Delta \to 0$.
\end{assumption}

In the zero-noise limit, the dynamically admissible set reduces to:
\[
\Gamma^0 = \left\{ [X] \in \mathcal{A} \;\middle|\;
\exists\, [X_0] \in \mathcal{A},\; \exists\, t \in [0,T] \text{ s.t. }
\Phi^t_F([X_0]) = [X] \right\}
\]
where $\Phi^t_F$ is the flow of $F$ at time $t$.

\begin{definition}[Admissible Tangent Cone]
For $[X] \in \Gamma^0$:
$\mathcal{T}_{[X]}\Gamma^0 = \{v \in T_{[X]}\mathcal{A} \mid v = F([X];\theta)
\text{ for some } \theta \in \Theta\}$.
\end{definition}

\begin{proposition}[Tangent Dimension Reduction]
If $\Theta \subset \mathbb{R}^p$, then $\dim \mathcal{T}_{[X]}\Gamma^0 \le p$.
In particular, if $p < \dim T_{[X]}\mathcal{A}$, the dynamical constraint reduces the
effective dimensionality of admissible directions from $\dim T_{[X]}\mathcal{A}$ to
at most $p$.
\end{proposition}
\begin{proof}
The map $\theta \mapsto F([X];\theta)$ is smooth from $\Theta \subset \mathbb{R}^p$ to
$T_{[X]}\mathcal{A}$; its image has dimension at most $p$ by the rank theorem. $\square$
\end{proof}

\subsection{Trajectory Learning from Before/After Pairs}

\begin{definition}[Fragment-Consistent Parameter Set]
Given fragment $(y_{\mathrm{before}}, y_{\mathrm{after}}, \tau)$:
\[
\Theta_{(y_b,y_a,\tau)} = \left\{ \theta \in \Theta \;\middle|\;
\exists\, [X_b] \in \Omega_{\mathrm{obs}}(y_b),\;
\Pi_i(\Phi^\tau_F([X_b])) \in B_{\epsilon_i}(y_{a,i})\; \forall i \right\}.
\]
Given training set $\mathcal{D}$, the learned parameter set is
$\hat{\Theta} = \bigcap_{n} \Theta_{(y_b^{(n)}, y_a^{(n)}, \tau^{(n)})}$.
\end{definition}

\textbf{The operator class library.} In practice $F(\cdot\,;\theta)$ is decomposed as:
\[
F([X];\theta) = \sum_{k=1}^K \alpha_k([X]) \cdot F_k([X];\theta_k)
\]
where each $F_k$ corresponds to a named ecological or civilizational process---forest
succession, urban expansion, agricultural cycling, infrastructure decay, hydrological
change, demographic transition. Each $F_k$ is learned from before/after pairs in which
that process is dominant. The activation weights $\alpha_k([X]) \ge 0$ satisfy
$\sum_k \alpha_k = 1$.

\subsection{The Dynamical Disambiguation Theorem}

\begin{definition}[Trajectory Separation]
Two configurations $[X],[X'] \in \Omega_{\mathrm{obs}}$ are \emph{trajectory-separated} by
$F(\cdot\,;\hat{\theta})$ at time $\tau > 0$ if
\[
\left\| \Ftil(\Phi^\tau_F([X];\hat{\theta})) - \Ftil(\Phi^\tau_F([X'];\hat{\theta}))
\right\|_{\prod_i \mathcal{Y}_i} > \epsilon
\]
for some $\epsilon > 0$ depending only on the noise levels $\{\epsilon_i\}$.
\end{definition}

\begin{theorem}[Dynamical Disambiguation]
\label{thm:disambiguation}
Under Assumptions~\ref{ass:regularity}--\ref{ass:generator}, if $|\Omega_{\mathrm{obs}}| > 1$
and every pair $[X],[X'] \in \Omega_{\mathrm{obs}}$ with $[X] \ne [X']$ is
trajectory-separated by $F(\cdot\,;\hat{\theta})$ at some $\tau \in (0,T]$, then
$|\mathcal{F}| = |\Omega_{\mathrm{obs}} \cap \Gamma^0| = 1$.
\end{theorem}

\begin{proof}
Let $[X],[X'] \in \Omega_{\mathrm{obs}} \cap \Gamma^0$ be distinct. Since both lie in
$\Gamma^0$, both lie on integral curves of $F(\cdot\,;\hat{\theta})$. By trajectory
separation, there exists $\tau \in (0,T]$ such that the forward-evolved projections differ
by more than $\epsilon$. At time $\tau$, observed data $y(\tau)$ must lie within $\epsilon_i$
of the true trajectory; therefore at most one of $[X(\tau)],[X'(\tau)]$ lies in
$\Omega_{\mathrm{obs}}(y(\tau))$. The true state $[X^*(\tau)] \in \Omega_{\mathrm{obs}}(y(\tau))$
by Lemma 5.1 at time $\tau$. Applying to all pairs eliminates all but $[X^*]$.
Hence $|\mathcal{F}| = 1$. $\square$
\end{proof}

\subsection{Joint Identifiability}

\begin{corollary}[Joint Identifiability]
The consistency operator $\mathcal{C}$ restricted to $\mathcal{F}$ has a unique fixed point
if either: (i) $\Ftil$ is injective on $\mathcal{F}$; (ii) every pair in
$\Omega_{\mathrm{obs}} \cap \Gamma^0$ is trajectory-separated; or (iii) both conditions
hold partially. The sufficient condition for (iii) is that $\dim(\Omega_{\mathrm{obs}}/\ker\Ftil)
+ \dim(\Gamma^0/\Gamma^0_{\mathrm{indist}}) \ge \dim\Omega_{\mathrm{obs}}$.
\end{corollary}

Case (iii) is the generic realistic scenario: sensors eliminate most ambiguity, trajectory
learning eliminates the residual. The two mechanisms are additive in their dimensional
reduction, not competing.

%% ============================================================
\section{The Inversion Layer}
\label{sec:inversion}
%% ============================================================

\preliminary{
\textbf{This section develops the inversion layer as an architecture for approximating the
consistency operator in practice, without conflating the approximation with the operator
itself.}

\medskip
\textbf{7.1 Role Specification.}
States the division of labor precisely: the inversion layer is not the model, it is a
proposal mechanism. Its role is to place the iterative optimization in the basin of the
correct fixed point. The consistency operator $\mathcal{C}$ is the model. The encoder is
a high-entropy initializer; the physics-constrained forward model is the contraction
operator. Explains why conflating encoder with solution is the failure mode of standard
encoder-decoder stacks: doing so makes the failure modes of Corollary~\ref{cor:failures}
undiagnosable. Separating proposal from decision preserves the falsifiability structure.

\medskip
\textbf{7.2 The Encoder as Approximate Posterior.}
Defines the encoder as implementing an approximate posterior:
\[
q_\phi(X \mid \{y_i\}) \approx p(X \mid \{y_i\}) \propto
\exp\!\left(-\sum_i d_i(\Pi_i(X), y_i) - \mathcal{R}(X)\right).
\]
The encoder approximates this distribution, not solves it. The forward-consistency loop
is what contracts that distribution. Discusses the variational inference interpretation:
the encoder minimizes KL divergence from the true posterior, subject to the architectural
constraints of the learned network. Notes that because $\mathcal{R}$ encodes physical
admissibility, the approximate posterior is already biased toward physically plausible
configurations---the encoder does not need to learn physics from scratch.

\medskip
\textbf{7.3 The Refinement Loop.}
Specifies the five-step loop:
\begin{enumerate}
\item Encoder samples $X^{(0)} \sim q_\phi(\cdot \mid \{y_i\})$
\item Forward model computes $\{\hat{y}_i\} = \{\Pi_i(X^{(0)})\}$
\item Residuals $\{r_i\} = \{d_i(\hat{y}_i, y_i)\}$ drive gradient update
\item $X^{(k+1)} = \mathcal{C}(X^{(k)})$ (one step of projected gradient descent on $f$)
\item Iterate until $\dim \mathcal{F} < \varepsilon$ (closure criterion from
  Section~\ref{sec:closure})
\end{enumerate}
Discusses convergence: Theorem~\ref{thm:identifiability} guarantees a unique fixed point
when conditions hold; the refinement loop approximates the path to that fixed point.
Convergence rate depends on the condition number of the Hessian of $f$ at $[X^*]$, which
is related to the Fisher information geometry of Appendix~\ref{app:fisher}.

\medskip
\textbf{7.4 The Categorical Structure of Multi-Modal Fusion.}
Develops the secondary, structural description: the encoder as a collection of
modality-specific proposal maps composed via a lax monoidal structure over
$\prod_i \mathcal{Y}_i$. Each modality-specific encoder $q_\phi^{(i)}$ produces a
marginal proposal; the fusion map combines them before the consistency operator acts.
This describes how multi-modal proposals are combined, but the explanatory weight rests
on the posterior formulation of Section 7.2, not on the categorical structure. The
categorical framing is reserved for structural analysis of how proposals compose across
scales in Section~\ref{sec:deployment}.

\medskip
\textbf{7.5 The Forward Model.}
Specifies what the forward model must compute: given a proposed $[X]$, produce
$\{\Pi_i(X)\}$ using the projection operators defined in Section~\ref{sec:projections}.
For infrared and LiDAR, this is direct evaluation of the operators. For RGB, it is a
volume rendering pass. For audio, it requires solving the RSVP field equations forward in
time from $X$ to produce $\partial_t S$ and $\nabla \cdot \mathbf{v}$ at microphone
locations. Discusses the computational cost of each projection and strategies for
approximation: learned surrogate models for the expensive projections (audio, photonic
strip), analytic projections for the cheap ones (infrared boundary traces).

\medskip
\textbf{7.6 Why Not End-to-End.}
Argues explicitly against treating the encoder as the final answer. An end-to-end system
collapses proposal and decision into a single learned map from $\{y_i\}$ to $[X]$.
This is equivalent to replacing $\mathcal{C}$ with the encoder and discarding the
forward model. The result is a system that cannot diagnose its own failure modes, cannot
be updated when new sensors are added (without retraining), and cannot be falsified by
physical constraints. The separation of encoder from consistency operator is not an
engineering choice---it is required by the identifiability structure.
}

%% ============================================================
\section{Closure: Definition, Criteria, and Detection}
\label{sec:closure}
%% ============================================================

\preliminary{
\textbf{This section defines closure operationally and proves that the three closure
criteria are equivalent under the main theorems.}

\medskip
\textbf{8.1 The Unified Definition.}
A state $[X^*] \in \Xmod$ is \emph{closed} if it is the unique fixed point of
$\mathcal{C}$ and lies on a trajectory of $T_\Delta$ consistent with all observed
fragments. This definition is then decomposed into three conditions---projection parity,
dynamical stability, counterfactual consistency---which are the operational handles on
closure in a deployed system.

\medskip
\textbf{8.2 The Three Conditions.}
\textit{Projection parity}: $\|\Pi_i(X^*) - y_i\|_{\mathcal{Y}_i} \le \epsilon_i$ for all
$i \in \mathcal{I}$. This is the basic requirement that the fixed point is consistent with
all current observations. It is verifiable in real time.

\textit{Dynamical stability}: $[X^*] \in \Gamma^0$, meaning $[X^*]$ lies on an integral
curve of $F(\cdot\,;\hat{\theta})$. Equivalently, the fixed point is not merely consistent
with instantaneous observations but is reachable under the learned dynamics. This rules
out configurations that happen to fit current observations but could not have been produced
by the world's actual trajectory.

\textit{Counterfactual consistency}: the system must be stable under re-projection. If
$[X^*]$ is simulated forward under $F$ for time $\tau$ and the resulting
$[\hat{X}^*(\tau)]$ is projected back to observation space, the result must lie within
$\epsilon_i$ of the actual observations at time $\tau$. This is the strongest condition:
it requires not just that the fixed point is consistent with past observations but that
its forward trajectory is consistent with future observations.

\medskip
\textbf{8.3 Equivalence Theorem.}
States and proves that the three conditions are equivalent under the assumptions of
Theorems~\ref{thm:identifiability} and \ref{thm:disambiguation}. The key step is showing
that each condition implies $\dim \mathcal{F} = 0$, and that $\dim \mathcal{F} = 0$
implies all three. The proof uses the identification from Section~\ref{sec:equivalence}:
the three conditions are projection parity ($\Omega_{\mathrm{obs}}$ tight), dynamical
stability ($\Gamma$ non-trivially intersecting $\Omega_{\mathrm{obs}}$), and counterfactual
consistency (the intersection is trajectory-stable, not just pointwise). These are three
faces of $|\mathcal{F}| = 1$.

\medskip
\textbf{8.4 Operational Detection.}
Closure is operationally detectable as stability of $[X^*]$ under perturbation of
$\epsilon_i$: if tightening any constraint by $\delta\epsilon_i$ does not move the fixed
point, closure has been achieved. This gives a practical stopping criterion for the
refinement loop of Section~\ref{sec:inversion}. Discusses the rate of convergence to
closure as a function of sensor density and trajectory fragment coverage, connecting to
the Fisher information geometry of Appendix~\ref{app:fisher}.

\medskip
\textbf{8.5 Closure and Entropy Collapse.}
Connects the closure detection criterion to the entropy identification of
Section~\ref{sec:equivalence}: at closure, $\vol_{\mathrm{Fisher}}(\mathcal{F}) = 0$,
$H^1(\mathcal{U},\mathcal{S}_\Pi) = 0$, $\dim H^0 = 1$, and $S_{\mathrm{RSVP}}(X^*)$
is minimized subject to admissibility---simultaneously and equivalently. The closed state
is not a best estimate; it is the unique trajectory consistent with all projections and
admissible dynamics. No alternative history remains consistent with the evidence.

\medskip
\textbf{8.6 Failure Mode Revisited.}
Revisits Corollary~\ref{cor:failures} in the light of the equivalence theorem. Projection
degeneracy ($H^1 \ne 0$) means closure cannot be achieved by any amount of additional
observation at the current moment---it requires either new sensor modalities or temporal
extension of the cover. Regularizer flatness means closure cannot be achieved by any
sensor---it requires revision of the model class $\mathcal{A}$. These are the only two
ways a Yarncrawler system can fail to converge, and both are diagnosable by the criteria
of this section.
}

%% ============================================================
\section{Deployment Scales as Feasible Set Geometry}
\label{sec:deployment}
%% ============================================================

\preliminary{
\textbf{This section unifies planetary, regional, and site-scale deployment as the same
system with different geometries of the feasible set.}

\medskip
\textbf{9.1 The Unified Framework.}
States the central claim: all three deployment scales are instances of the same
consistency operator $\mathcal{C}$ with different $(\Pi_i, \mathcal{N}_i, \epsilon_i)$
profiles. The scale is not a separate architectural choice---it is a parameter of the
noise model. A planetary-scale deployment has large $\epsilon_i$ (sparse, low-frequency
observations); a site-scale deployment has small $\epsilon_i$ (dense, high-frequency
observations). The feasible set $\Omega_{\mathrm{obs}}$ has different geometry at each
scale, but the identifiability theorem and disambiguation theorem apply uniformly.

\medskip
\textbf{9.2 Planetary Scale.}
Sensors: satellite multispectral imagery (RGB and infrared), occasional LiDAR passes,
no audio, no photonic strips. Temporal resolution: weeks to months. Spatial resolution:
10--100 meters per pixel. Noise levels $\epsilon_i$ are large. Consequence: $\Omega_{\mathrm{obs}}$
is high-dimensional and weakly constrained. The system operates in the high-entropy
proposal regime, relying heavily on $\mathcal{R}$ to restrict $\mathcal{F}$. Trajectory
learning carries most of the disambiguation burden: seasonal cycles, annual land cover
change, and multi-year succession patterns provide the operator class library entries that
contract $\Gamma$. States the modified identifiability condition for this regime: the
projection family need not be complete (Proposition~\ref{prop:completeness} fails at
planetary scale due to missing audio and strip modalities), but trajectory learning can
compensate via Corollary 6.4(iii).

\medskip
\textbf{9.3 Regional Scale.}
Sensors: aerial multispectral and LiDAR surveys, possibly distributed acoustic sensors,
periodic photonic strip transects. Temporal resolution: days to weeks. Spatial resolution:
1--10 meters. Noise levels intermediate. Consequence: $\Omega_{\mathrm{obs}}$ contracts
along spatial dimensions but remains temporally undersampled. The system requires both
static projection richness (more modalities available than at planetary scale) and
trajectory learning. The operator class library is expanded: finer-grained process
operators for agricultural change, urban growth, forest disturbance. Discusses the
transition from planetary to regional scale as a cover refinement: the planetary-scale
$\mathcal{S}_\Pi$ restricts to the regional-scale $\mathcal{S}_{\Pi'}$ on the finer
cover, with $H^1$ classes killed by the additional modalities.

\medskip
\textbf{9.4 Site Scale.}
Sensors: fixed arrays including all five modalities at high spatial and temporal density.
Temporal resolution: seconds to minutes. Spatial resolution: centimeters to meters. Noise
levels $\epsilon_i$ small. Consequence: $\Omega_{\mathrm{obs}}$ collapses rapidly toward
a singleton. The system achieves genuine closure and can generate counterfactual
re-projections for real-time validation. Discusses the calibration loop: at site scale,
re-projection residuals (forward model outputs vs.\ actual sensor readings) provide a
continuous closure detection signal, implementing the operational criterion of
Section~\ref{sec:closure}.

\medskip
\textbf{9.5 Multi-Scale Fusion.}
The three scales are not separate deployments---they are nested constraint systems. States
the Scale Consistency Theorem: if the noise models $\{\mathcal{N}_i\}$ are correctly
specified at each scale, the fixed point $[X^*]$ is invariant to which scale drives the
final refinement step. Proves this as a consequence of the functor $\Psi$: the
restriction map from the finer to the coarser sheaf preserves the global section, so the
fixed point at site scale restricts to the fixed point at regional scale and then to the
fixed point at planetary scale. In practice this means planetary constraints provide the
prior on $\mathcal{R}$, regional constraints initialize $q_\phi$, and site constraints
drive the refinement loop to closure.

\medskip
\textbf{9.6 Relaxing the Completeness Assumption.}
Returns to Proposition~\ref{prop:completeness} and asks what happens when the full
five-modality family is not available. Characterizes the residual ambiguity in
$\Omega_{\mathrm{obs}}$ as a function of which modalities are missing: missing audio
leaves $\partial_t S$ and $\nabla \cdot \mathbf{v}$ unconstrained (the entropy production
and flow divergence degrees of freedom remain free); missing LiDAR leaves the level-set
geometry of $\Phi$ ambiguous; and so on. Each missing modality corresponds to a direction
in $\ker \Ftil|_{\mathcal{F}}$ that can only be resolved by trajectory learning. This
gives a sensor selection criterion grounded in the geometry of the feasible set: the most
valuable next sensor to add is the one that closes the largest-dimensional remaining
direction in $\ker \Ftil|_{\mathcal{F}}$.
}

%% ============================================================
\section{Relation to Prior Work}
\label{sec:related}
%% ============================================================

\preliminary{
\textbf{This section positions Yarncrawler against four traditions: data assimilation,
inverse rendering, digital twin systems, and categorical/sheaf-theoretic semantics.
The goal is not to survey those traditions but to identify precisely where Yarncrawler
departs from each and why the departure is necessary.}

\medskip
\textbf{10.1 Data Assimilation and Ensemble Methods.}
Data assimilation---especially ensemble Kalman filtering (EnKF) and variational methods
(4D-Var)---solves the same broad problem: estimating a hidden state from observations
over time. The key differences are: (a) data assimilation works in $\mathcal{X}$ rather
than $\Xmod$, treating gauge redundancies as noise rather than structure; (b) the
regularizer in data assimilation is a prior distribution, not an admissibility condition---
it expresses preference, not exclusion; (c) data assimilation has no sheaf structure:
there is no notion of local charts, restriction maps, or cohomological obstruction. The
Yarncrawler Identifiability Theorem has no counterpart in standard data assimilation
because the quotient structure is absent. Notes where EnKF and 4D-Var are special cases:
in the linear Gaussian regime with trivial symmetry group ($G = \{e\}$), $\mathcal{C}$
reduces to the Kalman update.

\medskip
\textbf{10.2 Inverse Rendering and Neural Radiance Fields.}
Neural radiance fields (NeRF) and related inverse rendering methods solve the problem of
reconstructing a 3D scene from 2D images. The RGB projection operator $\Pi_{\mathrm{RGB}}$
defined in Section~\ref{sec:projections} is the NeRF rendering integral. The differences
are: (a) NeRF has no dynamics---it is a purely static reconstruction method with no
$\Gamma$; (b) NeRF has no RSVP field structure---the scene representation is a learned
radiance function, not a physical field triple; (c) NeRF has no regularizer encoding
physical admissibility. The Yarncrawler framework strictly generalizes NeRF: the RGB
projection is one chart of the cover, and NeRF is the degenerate case where only that
chart is available and dynamics are absent.

\medskip
\textbf{10.3 Digital Twin Systems.}
Digital twin literature proposes maintaining synchronized virtual replicas of physical
systems. The standard approach is object-ontological: the twin represents a collection of
objects (buildings, trees, vehicles) whose properties are updated from sensor data.
Yarncrawler departs in three ways: (a) the primary ontology is trajectory-based, not
object-based---objects emerge as persistent patterns in field trajectories, not as primary
entities; (b) the twin is not updated by recognized objects but by the consistency
operator acting on the full quotient configuration space; (c) digital twins have no
identifiability theory---there is no formal account of when a twin uniquely determines
its physical counterpart. The Yarncrawler Identifiability Theorem provides exactly this.

\medskip
\textbf{10.4 Categorical and Sheaf-Theoretic Semantics.}
Sheaf theory has been applied to contextuality in quantum mechanics, to distributed
knowledge structures in multi-agent systems, and to compositional semantics in natural
language processing. Yarncrawler draws on these applications but departs in one crucial
direction: the sheaf here is over a physical configuration space with RSVP field structure,
not over an abstract semantic space. The stalks are not sets of propositions or probability
distributions but Sobolev-class field configurations satisfying physical equations. The
repair morphisms are not logical updates but steps of a variational optimization. The
entropy identified with $H^1$ is a physical entropy field, not an information-theoretic
measure. This grounding is what allows the framework to make contact with sensor physics
and deployment engineering while retaining the full generality of the categorical structure.

\medskip
\textbf{10.5 The Yarncrawler Program.}
Situates the present paper within the broader Yarncrawler program, which applies the same
sheaf-repair-over-RSVP architecture to semantic systems, cultural evolution, ecological
dynamics, and AI reasoning. The present paper is the inverse direction of that program:
where the forward direction takes a state and asks how coherence is maintained under
perturbation, the inverse direction takes observations and asks what state they demand.
The Sheaf-Variational Equivalence of Section~\ref{sec:equivalence} is the theorem that
makes the two directions formally dual rather than merely analogous.
}

%% ============================================================
\section{Discussion: Closure as Epistemology}
\label{sec:epistemology}
%% ============================================================

\preliminary{
\textbf{This section draws the epistemological consequence of the framework and connects
it to the broader Yarncrawler theoretical program.}

\medskip
\textbf{11.1 What Closure Means.}
A closed Yarncrawler state is not a probabilistic estimate. It is not the most likely
world-state given the observations. It is the unique world-state that is structurally
demanded by the observations under the model class. The distinction is not merely
philosophical: it has operational consequences. A probabilistic estimate can always be
revised by new observations; a closed state cannot be revised without either adding new
sensor modalities (killing new $H^1$ classes) or revising the model class (changing
$\mathcal{A}$). The closed state is, in a precise sense, as certain as the model class
allows.

\medskip
\textbf{11.2 Constraint Primacy.}
The framework operationalizes a philosophical commitment: constraints are prior to
completions. The world-state is not recovered by filling in a model template with
observations---it is the unique configuration that survives after all constraints have
been applied. This is the physical analogue of a principle that recurs throughout the
Yarncrawler program: meaning emerges from constraint closure rather than from
representation. The Sheaf-Variational Equivalence makes this precise: a world-state has
meaning (is uniquely identified) exactly when the sheaf induced by its projections has a
unique global section.

\medskip
\textbf{11.3 The RSVP Fields as Constraint Carriers.}
The RSVP field triple $(\Phi, \mathbf{v}, S)$ is not merely a convenient parametrization
of the configuration space. It is the structure that makes the regularizer $\mathcal{R}$
non-arbitrary: the admissible set $\mathcal{A}$ is defined by the RSVP field equations,
which encode physical conservation laws, ecological succession principles, and
civilizational feasibility constraints. Without the RSVP structure, the regularizer would
be an arbitrary penalty function and the closed state would be an artifact of the choice
of penalty. With it, the closed state is the unique configuration consistent with the laws
governing the domain. The framework is therefore not just an inference procedure---it is
a theory of what it means for a physical configuration to be real.

\medskip
\textbf{11.4 Stigmergic Accumulation and the Refinement Loop.}
Connects the refinement loop of Section~\ref{sec:inversion} to the stigmergic dynamics
developed in the Yarncrawler forward direction. Each step of the refinement loop is a
repair morphism in the sheaf $\mathcal{S}_\Pi$: it reduces a projection residual
(trivializes a cocycle) and updates the field configuration toward closure. The
accumulation of constraints over time---as more sensor observations arrive---is the
variational analogue of stigmergic reinforcement: each new observation narrows the
feasible set, and the trajectory of the refinement loop through configuration space is
the path the crawler takes through the yarn. The loop converges when no more constraints
can be satisfied---when the yarn has been fully crawled and the global section is unique.

\medskip
\textbf{11.5 Limits of the Framework.}
Acknowledges three structural limits:

\textit{Model class dependence}: the closed state is unique within $\mathcal{A}$, not in
$\mathcal{X}$ or $\Xmod$. If $\mathcal{A}$ is misspecified---if the true world-state lies
outside the admissible set---the consistency operator will converge to the closest
admissible state, not the true state. Non-existence of a global section (discussed in
Section~\ref{sec:closure}) is the signal that model revision is needed.

\textit{Trajectory learning quality}: the dynamical disambiguation result depends on the
quality of $F(\cdot\,;\hat{\theta})$. If the learned dynamics are systematically wrong---
if the operator class library is missing a dominant process type---trajectory separation
will fail for pairs of configurations that differ primarily along the missing process
direction. This is detectable as residual $|\mathcal{F}| > 1$ after full sensor coverage.

\textit{Computational tractability}: the consistency operator is defined as a global
minimizer of a non-convex functional in infinite dimensions. The refinement loop
approximates this with local gradient steps. Convergence to the global minimizer is
guaranteed by the identifiability theorem only up to initialization quality; in practice
multiple initializations from $q_\phi$ may be needed.

\medskip
\textbf{11.6 Forward and Inverse as Dual Inference.}
Closes by returning to the duality of Section 1.3. The forward direction maintains
coherence by repairing tears as they arise: it is the epistemic activity of a system
embedded in a world it partially controls. The inverse direction reconstructs state from
observations: it is the epistemic activity of a system observing a world it does not
control. Most real intelligent systems do both simultaneously. The Yarncrawler framework
is the unified theory of both activities---the sheaf $\mathcal{S}_\Pi$ is the common
object, and the two directions are two ways of relating to it. A fully realized Yarncrawler
system would alternate between the two: reconstructing the world-state from observations
(inverse direction), projecting forward to predict future states (forward direction),
comparing predictions to observations to detect model failure (closure detection), and
repairing the model when failure is detected (sheaf repair). This is the complete
epistemic loop of a system that both observes and maintains.
}

%% ============================================================
\section{Conclusion}
\label{sec:conclusion}
%% ============================================================

\preliminary{
\textbf{The conclusion is short---three to four paragraphs. It does not summarize the
paper. It states what the paper has established, what remains to be done, and what the
framework implies.}

\medskip
\textbf{Paragraph 1: What has been established.}
The Yarncrawler Identifiability Theorem and the Dynamical Disambiguation Theorem together
characterize the conditions under which multi-modal sensor observations uniquely determine
a physical world-state. The Sheaf-Variational Equivalence shows that this characterization
is not merely a technical result about inverse problems: it is the variational realization
of a categorical structure in which sensor modalities are charts, observations are local
sections, and closure is the existence of a unique global section. The RSVP field triple
is the substrate in which these structures are grounded.

\medskip
\textbf{Paragraph 2: The main finding compressed.}
A world-state is uniquely recoverable from multi-modal observations if and only if the
induced map from the quotient configuration space to the product of observation spaces is
injective on the feasible set, or if the dynamics of the world separate configurations
that static sensors cannot. In either case, closure corresponds to cohomology
trivialization, entropy collapse, and Fisher volume collapse---simultaneously and
equivalently.

\medskip
\textbf{Paragraph 3: What remains.}
Three open problems are identified. (a) The extension of the identifiability theorem to
the stochastic case without the low-noise limit: does a version of the theorem hold when
$\sigma^2$ is not small? (b) The learning theory of $F(\cdot\,;\hat{\theta})$: what
sample complexity is needed for the operator class library to ensure trajectory separation
with high probability? (c) The connection between the Fisher information geometry of the
feasible set and the computational complexity of the refinement loop: can the convergence
rate of the refinement loop be bounded in terms of the eigenspectrum of the Fisher metric
at $[X^*]$?

\medskip
\textbf{Paragraph 4: The implication.}
The Yarncrawler system is not a sensing architecture. It is a theory of when the world
becomes legible---when the observations accumulated by a sensor network, constrained by
physical law and dynamical history, leave no alternative configuration standing. That
condition is not an approximation target. It is a threshold: on one side, multiple
histories remain consistent with the evidence; on the other, only one does. The framework
identifies that threshold precisely and provides the operator that reaches it.
}

%% ============================================================
%% APPENDICES
%% ============================================================

\appendix

%% ============================================================
\section{Fréchet Manifold Structure of $\mathcal{X}$}
\label{app:frechet}
%% ============================================================

\preliminary{
\textbf{This appendix establishes the differential-geometric foundations required for the
main theorems.}

\medskip
\textbf{A.1 Fréchet Manifolds.}
Defines Fréchet manifolds as locally convex topological vector spaces with smooth
transition maps. Notes that $\mathcal{X} = H^1(\Omega) \times H^1(\Omega;\mathbb{R}^3)
\times L^2(\Omega)$ is in fact a Hilbert space (hence also a Banach space and a Fréchet
space) with its product inner product. The Fréchet framing is introduced because
extensions of the framework to more general field configurations (e.g., distributional
fields, or fields on non-compact domains) may require Fréchet rather than Hilbert
structure.

\medskip
\textbf{A.2 The Quotient Topology.}
Proves that $\Xmod$ is Hausdorff under the assumption that $G$ acts properly and
isometrically. Gives the metric on $\Xmod$ induced from $\mathcal{X}$. Discusses the
difference between $\Xmod$ (observational equivalence) and $\mathcal{X}/G$
(gauge equivalence) and when they coincide.

\medskip
\textbf{A.3 Tangent Bundle of $\Xmod$.}
Defines the tangent bundle $T\Xmod$ as the quotient of $T\mathcal{X}$ by the action of
$TG$. Specifies the admissible tangent cone $\mathcal{T}_{[X]}\Gamma^0$ (Definition 6.7
of the main text) in terms of this bundle structure. Notes that when $G$ acts freely,
$T\Xmod \cong T\mathcal{X}/TG$ is a vector bundle, and the tangent cone is a linear
subspace.

\medskip
\textbf{A.4 Geodesic Convexity on $\Xmod$.}
Defines geodesics on $\Xmod$ via the quotient of geodesics on $\mathcal{X}$ (which are
straight lines in the Hilbert structure). Proves that $\mathcal{A}$ is geodesically convex
when its preimage in $\mathcal{X}$ is convex---which holds by the definition of
$\mathcal{A}$ as the zero-sublevel set of a convex functional $\mathcal{R}$. This justifies
Assumption A4 (Assumption~\ref{ass:admissibility} in the main text).
}

%% ============================================================
\section{Full Proof of the Identifiability Theorem}
\label{app:identifiability}
%% ============================================================

\subsection{Functional-Analytic Setting}

The proof of Theorem~\ref{thm:identifiability} relies on existence and uniqueness of
minimizers of the consistency functional on $\Xmod$. We introduce additional assumptions
for the infinite-dimensional setting.

\begin{assumption}[Banach Realization, restated]
There exists a Banach space $\mathcal{B}$ and continuous embedding
$\iota : \mathcal{X} \hookrightarrow \mathcal{B}$ such that $\mathcal{A}$ is closed and
convex in $\mathcal{B}$, each $\Pi_i$ extends to a bounded linear operator on $\mathcal{B}$,
and $\mathcal{R}$ extends to a proper, lower semicontinuous functional on $\mathcal{B}$.
In practice, $\mathcal{B} = \mathcal{X}$ itself.
\end{assumption}

\begin{assumption}[Geodesic Convexity, restated]
$\mathcal{A}$ is geodesically convex with respect to the metric induced from $\mathcal{B}$,
and $\mathcal{R}$ is strictly convex along geodesics in $\mathcal{A}$:
\[
\mathcal{R}(\gamma(t)) < (1-t)\mathcal{R}([X]) + t\mathcal{R}([X'])
\quad \forall [X] \ne [X'] \in \mathcal{A},\; t \in (0,1).
\]
\end{assumption}

\begin{assumption}[Coercivity]
$f([X]) = \sum_i \|\Pi_i(X)-y_i\|^2 + \mathcal{R}(X)$ satisfies
$\|X_n\|_{\mathcal{B}} \to \infty \Rightarrow f([X_n]) \to \infty$.
\end{assumption}

\subsection{Existence of a Minimizer}

\begin{proposition}
Under the standing assumptions, $f$ attains a minimum on $\mathcal{A}$.
\end{proposition}
\begin{proof}
Let $\{[X_n]\}$ be a minimizing sequence; coercivity gives boundedness in $\mathcal{B}$.
Reflexivity of $\mathcal{B}$ (if Hilbert) yields a weakly convergent subsequence
$X_{n_k} \rightharpoonup X_*$. Closedness of $\mathcal{A}$ gives $[X_*] \in \mathcal{A}$.
Weak continuity of $\Pi_i$ (bounded linear) and lower semicontinuity of $\mathcal{R}$
give $f([X_*]) \le \liminf f([X_{n_k}])$. Hence $[X_*]$ is a minimizer. $\square$
\end{proof}

\subsection{Uniqueness of the Minimizer}

\begin{proposition}
Under the standing assumptions, the minimizer is unique.
\end{proposition}
\begin{proof}
Suppose $[X],[X']$ are distinct minimizers. By geodesic convexity, there exists a
geodesic $\gamma$ in $\mathcal{A}$ connecting them. The projection residual term is
convex along $\gamma$ (by linearity of $\Pi_i$). The regularizer is strictly convex along
$\gamma$ by assumption. Hence $f(\gamma(t)) < (1-t)f([X]) + tf([X']) = \inf f$ for
$t \in (0,1)$, contradicting minimality. $\square$
\end{proof}

\subsection{Passage to the Quotient}

\begin{proposition}
If $[X] = [X']$ in $\Xmod$, then $f([X]) = f([X'])$. Conversely, distinct equivalence
classes yield distinct objective values whenever $\Ftil$ is injective.
\end{proposition}
\begin{proof}
Observational equivalence means $\Pi_i(X) = \Pi_i(X')$ for all $i$, so residual terms
are equal; $\mathcal{R}$ is defined on equivalence classes, so regularizer terms are
equal. Injectivity of $\Ftil$ means distinct classes differ in at least one projection,
so residual terms differ. $\square$
\end{proof}

\subsection{Full Proof of Theorem~\ref{thm:identifiability}}

\begin{proof}[Full proof]
\textit{Well-definedness}: by the existence and uniqueness propositions, $\mathcal{C}$ is
single-valued on $\mathcal{A}$.

\textit{Fixed points = minimizers}: $[X^*] = \mathcal{C}([X^*])$ iff $[X^*]$ minimizes
$f$ over $\mathcal{A}$.

\textit{Injectivity implies uniqueness}: if $\Ftil$ is injective on $\mathcal{F}$,
$\mathcal{F}$ contains at most one class. By non-emptiness, $[X^*] \in \mathcal{F}$, so
$|\mathcal{F}| = 1$. At $[X^*]$, all residuals are within tolerance; strict convexity
prevents alternative minimizers. Hence $[X^*]$ is the unique fixed point.

\textit{Non-injectivity implies non-uniqueness}: if $\Ftil$ is not injective on
$\mathcal{F}$, distinct $[X],[X'] \in \mathcal{F}$ have identical projections. Both are
candidates for the minimizer. If $\mathcal{R}([X]) = \mathcal{R}([X'])$, strict
convexity fails (contradiction). If $\mathcal{R}([X]) \ne \mathcal{R}([X'])$, the
larger-regularizer point is not a minimizer---but then it is not in $\mathcal{F}$ as
defined, contradicting the assumption. Either way we reach a contradiction with
uniqueness. $\square$
\end{proof}

%% ============================================================
\section{Full Proof of the Dynamical Disambiguation Theorem}
\label{app:disambiguation}
%% ============================================================

\subsection{Measure-Theoretic Framework}

\begin{assumption}[Polish Structure]
$\mathcal{A}$ with the metric induced from $\mathcal{B}$ is a Polish space (complete,
separable metric).
\end{assumption}

\begin{definition}[Markov Kernel, precise]
$T_\Delta : \mathcal{A} \to \mathcal{P}(\mathcal{A})$ is a Markov kernel: for each
$[X]$, $T_\Delta([X],\cdot)$ is a Borel probability measure on $\mathcal{A}$; for each
Borel $B \subseteq \mathcal{A}$, $[X] \mapsto T_\Delta([X],B)$ is measurable.
\end{definition}

\subsection{Convergence to the Deterministic Generator}

\begin{assumption}[Generator Convergence, precise]
For all bounded continuous $\varphi : \mathcal{A} \to \mathbb{R}$:
\[
\int \varphi(X')\,T_\Delta([X],dX') = \varphi([X] + \Delta F([X];\theta)) + O(\Delta^{3/2}).
\]
\end{assumption}

This is weak convergence $T_\Delta([X],\cdot) \Rightarrow \delta_{[X]+\Delta F([X];\theta)}$
at rate $\Delta^{3/2}$, corresponding to a diffusion with noise variance $O(\Delta)$.

\subsection{Stability of Trajectory Separation Under Noise}

\begin{definition}[Uniform Separation Margin]
A pair $[X],[X'] \in \Omega_{\mathrm{obs}}$ is \emph{uniformly trajectory-separated} if
there exist $\tau > 0$ and $\delta > 0$ such that
\[
\left\| \Ftil(\Phi^\tau_F([X])) - \Ftil(\Phi^\tau_F([X'])) \right\| \ge \epsilon + \delta.
\]
\end{definition}

\begin{lemma}
Under generator convergence, uniformly trajectory-separated pairs remain distinguishable
under stochastic evolution with probability $\ge 1 - C\Delta/\delta^2$.
\end{lemma}
\begin{proof}
By Chebyshev applied to the $O(\Delta)$ noise variance, stochastic trajectories deviate
from deterministic by $O(\sqrt{\Delta})$ with high probability. For deviation $\delta/2$:
$\mathbb{P}(\|X_\tau^\Delta - X_\tau\| > \delta/2) \le C\Delta/\delta^2$.
With margin $\delta$, the two stochastic trajectories remain separated by $\ge \epsilon$
with probability $\ge 1 - 2C\Delta/\delta^2$. $\square$
\end{proof}

\subsection{Full Proof of Theorem~\ref{thm:disambiguation}}

\begin{proof}[Full proof]
Let $[X],[X'] \in \Omega_{\mathrm{obs}} \cap \Gamma^0$ be distinct and uniformly
trajectory-separated at time $\tau$.

\textit{Step 1}: Both lie on integral curves of $F(\cdot\,;\hat{\theta})$ by definition of
$\Gamma^0$.

\textit{Step 2}: At time $\tau$, their projections differ by $\ge \epsilon + \delta > \epsilon$.

\textit{Step 3}: By the noise lemma, stochastic evolutions remain separated with high
probability for small $\Delta$.

\textit{Step 4}: At time $\tau$, the true observation $y(\tau)$ has
$\|y_i(\tau) - \Pi_i(X^*(\tau))\| \le \epsilon_i$. At most one of the two separated
candidates lies within $\epsilon_i$ of $y_i(\tau)$.

\textit{Step 5}: The true state $[X^*(\tau)] \in \Omega_{\mathrm{obs}}(y(\tau))$ by
non-emptiness applied at time $\tau$. Hence the false candidate is eliminated.

\textit{Step 6}: Applying to all pairs eliminates all but $[X^*]$. Hence
$|\mathcal{F}| = 1$. $\square$
\end{proof}

%% ============================================================
\section{Fisher Information Geometry of the Feasible Set}
\label{app:fisher}
%% ============================================================

\preliminary{
\textbf{This appendix develops the Riemannian geometry of $\mathcal{F}$ induced by the
noise models $\{\mathcal{N}_i\}$ and connects it to the entropy identification of
Section~\ref{sec:equivalence}.}

\medskip
\textbf{D.1 The Fisher Information Metric.}
Defines the Fisher information metric on $\Xmod$ induced by the family of observation
distributions $\{p_i(\cdot \mid [X]) = \mathcal{N}(\Pi_i(X), \Sigma_i)\}_{i \in \mathcal{I}}$.
For Gaussian noise models with covariance $\Sigma_i$, the Fisher metric at $[X]$ is:
\[
g_F([X])(u, v) = \sum_i \langle D\Pi_i(u), \Sigma_i^{-1} D\Pi_i(v) \rangle_{\mathcal{Y}_i}
\]
where $D\Pi_i$ is the derivative of $\Pi_i$ at $[X]$. This is a semi-Riemannian metric
on $\Xmod$: it may be degenerate in directions where all projections have zero derivative.

\medskip
\textbf{D.2 Volume of the Feasible Set.}
Defines $\vol_{\mathrm{Fisher}}(\mathcal{F})$ as the Riemannian volume of $\mathcal{F}$
under $g_F$. Shows that $\vol_{\mathrm{Fisher}}(\mathcal{F}) \propto \prod_i \epsilon_i^{n_i}$
where $n_i$ is the number of degrees of freedom closed by $\Pi_i$ not already closed by
$\{\Pi_1,\ldots,\Pi_{i-1}\}$ (the marginal contribution of each sensor). As $\epsilon_i \to 0$
for all $i$, $\vol_{\mathrm{Fisher}}(\mathcal{F}) \to 0$, corresponding to closure.

\medskip
\textbf{D.3 Connection to Entropy.}
Proves the identification
$\log \vol_{\mathrm{Fisher}}(\mathcal{F}) \leftrightarrow \log \dim H^0(\mathcal{U},\mathcal{S}_\Pi)
\leftrightarrow S_{\mathrm{RSVP}}(X^*)$
stated in Section~\ref{sec:equivalence}. The first identification uses the fact that
$\dim H^0$ counts linearly independent global sections, which in the Gaussian regime
scales with the volume of the feasible set. The second uses the RSVP entropy field
equation and the fact that at the fixed point $[X^*]$, $S_{\mathrm{RSVP}}$ is the
integral of local entropy production over $\Omega$, which is proportional to the
log-volume of remaining uncertainty.

\medskip
\textbf{D.4 Reconstruction Efficiency.}
Defines the \emph{reconstruction efficiency} of a sensor network as the rate of
$\vol_{\mathrm{Fisher}}(\mathcal{F})$ decrease per unit of added sensing capacity
(measured in total Fisher information $\sum_i \|\Sigma_i^{-1}\|_{\mathrm{op}}$). This
gives a principled criterion for sensor placement: the next sensor should be placed where
it maximizes the decrease in $\vol_{\mathrm{Fisher}}(\mathcal{F})$, which is determined
by the spectrum of the Fisher metric in the remaining unclosed directions.

\medskip
\textbf{D.5 Convergence Rate of the Refinement Loop.}
Bounds the convergence rate of the refinement loop (Section~\ref{sec:inversion}) in terms
of the eigenspectrum of $g_F$ at $[X^*]$. The refinement loop is a gradient descent on $f$;
its convergence rate is determined by the condition number $\kappa = \lambda_{\max}/\lambda_{\min}$
of the Hessian of $f$ at $[X^*]$, which is related to the eigenvalues of $g_F$ at $[X^*]$.
In particular, directions with small Fisher information (loosely constrained by sensors)
correspond to directions with small Hessian eigenvalues and hence slow convergence. This
gives a formal account of why dense, multi-modal sensor networks accelerate convergence
to closure.
}

%% ============================================================
\section{Operator Class Library}
\label{app:operators}
%% ============================================================

\preliminary{
\textbf{This appendix formally specifies the named transition operators in the operator
class library introduced in Section~\ref{sec:disambiguation}.}

\medskip
\textbf{E.1 Structure of the Library.}
The library consists of $K$ named process generators $\{F_k\}_{k=1}^K$, each a learned
vector field on $\mathcal{A}$ parametrized by $\theta_k \in \Theta_k$. The full dynamics
are $F([X];\theta) = \sum_k \alpha_k([X]) F_k([X];\theta_k)$ with
$\alpha_k \ge 0$, $\sum_k \alpha_k = 1$. Each $F_k$ is learned from before/after pairs
in which process $k$ is dominant, identified by its characteristic sensor signature.

\medskip
\textbf{E.2 Ecological Process Operators.}
Specifies the following operators with their characteristic sensor signatures:

\textit{Forest succession} ($F_{\mathrm{succ}}$): slow increase in LiDAR canopy height,
gradual shift in RGB reflectance from grass-like to canopy-like, decrease in audio
spectral entropy (bird diversity changes with succession stage). Time scale: years to
decades.

\textit{Disturbance and recovery} ($F_{\mathrm{dist}}$): rapid decrease in LiDAR height
and RGB NDVI, followed by slow recovery. Audio signature: absence of characteristic
species. Time scale: event-driven (days), recovery (years).

\textit{Hydrological dynamics} ($F_{\mathrm{hydro}}$): correlated changes in infrared
(soil moisture affects surface temperature), LiDAR (inundation changes geometry), and
audio (water flow sounds). Time scale: seasonal to event-driven.

\textit{Phenological cycling} ($F_{\mathrm{pheno}}$): regular seasonal oscillations in
RGB (green-up/senescence), infrared (latent heat flux changes), and audio (bird song
seasonality). Time scale: annual.

\medskip
\textbf{E.3 Civilizational Process Operators.}
Specifies:

\textit{Urban expansion} ($F_{\mathrm{urban}}$): progressive increase in LiDAR building
density and height, change in RGB to impervious surface reflectance, increase in audio
(traffic, machinery). Time scale: years to decades.

\textit{Infrastructure decay} ($F_{\mathrm{decay}}$): gradual change in LiDAR surface
roughness, spectral shift in RGB, increase in audio entropy (structural noise). Time
scale: years.

\textit{Agricultural cycling} ($F_{\mathrm{agri}}$): regular oscillations in RGB
(crop growth stages), infrared (ET changes), LiDAR (canopy height variation), audio
(machinery activity). Time scale: seasonal.

\textit{Demographic change} ($F_{\mathrm{demo}}$): correlated changes in audio activity
levels, RGB night-light proxies, and infrastructure use patterns. Time scale: years.

\medskip
\textbf{E.4 Learning Protocol.}
Specifies how each $F_k$ is learned from before/after pairs: (a) identify pairs where
process $k$ is the dominant change by classifying the observation delta against the
library; (b) fit $\theta_k$ by minimizing the discrepancy between the predicted
forward image $\Pi_i(\Phi^\tau_{F_k}([X_b]))$ and the actual after-observation $y_a$
over training pairs; (c) validate by held-out pairs and by checking that
$\Theta_{(y_b,y_a,\tau)} \cap \Theta_k$ is non-empty for pairs labeled as process $k$.

\medskip
\textbf{E.5 Library Extensibility.}
The library is open: new process types $F_{K+1}$ can be added by collecting before/after
pairs exhibiting the new process, fitting the generator, and appending to the library.
The identifiability and disambiguation theorems are not affected by library extension:
they hold for any $F$ with Lipschitz generator, regardless of how many process components
it contains. The effect of library extension is to reduce $\Gamma^0_{\mathrm{indist}}$
(the set of configurations not separated by trajectory learning), which by
Corollary 6.4(iii) increases the rate of closure.
}

\newpage
\bibliographystyle{plainnat}
\begin{thebibliography}{99}

\bibitem{friston2010}
Friston, K. (2010).
The free-energy principle: a unified brain theory?
\textit{Nature Reviews Neuroscience}, 11(2), 127--138.

\bibitem{friston2019}
Friston, K., FitzGerald, T., Rigoli, F., Schwartenbeck, P., \& Pezzulo, G. (2019).
Active inference: a process theory.
\textit{Neural Computation}, 31(1), 1--49.

\bibitem{hordijk2019}
Hordijk, W., \& Steel, M. (2019).
Autocatalytic networks: an intimate relation between network topology and dynamics.
\textit{Bulletin of Mathematical Biology}, 81(5), 1357--1374.

\bibitem{maturana1980}
Maturana, H. R., \& Varela, F. J. (1980).
\textit{Autopoiesis and Cognition: The Realization of the Living}.
D.~Reidel.

\bibitem{rosen1985}
Rosen, R. (1985).
\textit{Anticipatory Systems}.
Pergamon Press.

\bibitem{abramsky2011}
Abramsky, S., \& Brandenburger, A. (2011).
The sheaf-theoretic structure of non-locality and contextuality.
\textit{New Journal of Physics}, 13, 113036.

\bibitem{spivak2014}
Spivak, D. I. (2014).
\textit{Category Theory for the Sciences}.
MIT Press.

\bibitem{lewis2020}
Lewis, P., Perez, E., et al. (2020).
Retrieval-augmented generation for knowledge-intensive NLP tasks.
\textit{Advances in Neural Information Processing Systems}, 33.

\bibitem{mildenhall2021}
Mildenhall, B., Srinivasan, P. P., Tancik, M., Barron, J. T., Ramamoorthi, R.,
\& Ng, R. (2021).
NeRF: representing scenes as neural radiance fields for view synthesis.
\textit{Communications of the ACM}, 65(1), 99--106.

\bibitem{freidlin1984}
Freidlin, M. I., \& Wentzell, A. D. (1984).
\textit{Random Perturbations of Dynamical Systems}.
Springer.

\end{thebibliography}

\end{document}
